\documentclass[11pt,a4paper]{article}

% ---------- Codificación y tipografía ----------
\usepackage[utf8]{inputenc}
\usepackage[T1]{fontenc}
\usepackage{lmodern}
\usepackage[spanish,es-nodecimaldot]{babel}
\usepackage{microtype}
\usepackage{amsmath,amssymb}
\usepackage{geometry}
\geometry{top=2.3cm,bottom=2.3cm,left=2.2cm,right=2.2cm}
\usepackage{enumitem}
\usepackage{booktabs}
\usepackage{longtable}
\usepackage{array}
\usepackage{fancyhdr}
\usepackage{titlesec}
\usepackage{setspace}
\usepackage{multicol}
\usepackage{needspace}

\setstretch{1.08}
\setlength{\parindent}{0pt}
\setlength{\parskip}{4pt}

% ---------- Encabezado y pie ----------
\pagestyle{fancy}
\fancyhf{}
\lhead{\small Banco de Preguntas PEU --- SUNASS}
\rhead{\small Módulo X: Finanzas}
\cfoot{\small \thepage}
\renewcommand{\headrulewidth}{0.4pt}
\renewcommand{\footrulewidth}{0pt}

% ---------- Títulos ----------
\titleformat{\section}{\large\bfseries}{\thesection}{0.8em}{}
\titlespacing*{\section}{0pt}{14pt}{6pt}
\titleformat{\subsection}{\normalsize\bfseries}{\thesubsection}{0.6em}{}
\titlespacing*{\subsection}{0pt}{10pt}{4pt}

% ---------- Comandos del examen ----------
\newcounter{qnum}
\newcommand{\q}[1]{%
  \needspace{8\baselineskip}%
  \par\medskip\refstepcounter{qnum}%
  \noindent\textbf{\theqnum.}\ #1\par\nopagebreak[4]
}
\newenvironment{opts}
  {\begin{enumerate}[label=\Alph*),leftmargin=2.0em,itemsep=1pt,topsep=2pt,parsep=0pt,partopsep=0pt]}
  {\end{enumerate}}
\makeatletter
\@itempenalty=10000
\makeatother
\newcommand{\sol}[3]{\par\noindent\textbf{#1. #2.}\ #3\par\smallskip}

% ==============================================================================
% Configuración BackBone Academy para los Bancos de Preguntas (Módulos 01-10)
% ==============================================================================
\usepackage{graphicx}
\usepackage{xcolor}
\usepackage{tikz}
\usepackage[hidelinks]{hyperref}

% ==============================================================================
% Identidad de marca BackBone Academy compartida (Beamer y Bancos de Preguntas)
% ==============================================================================
\definecolor{azulPrincipal}{RGB}{43, 45, 162}
\definecolor{azulOscuro}{RGB}{24, 26, 98}
\definecolor{grisTexto}{RGB}{85, 85, 95}

\newcommand{\backboneIgUrl}{https://www.instagram.com/backbone.academy}
\newcommand{\backboneTtUrl}{https://www.tiktok.com/@backbone.academy}
\newcommand{\backboneWebUrl}{https://leonardovelasquez14-detg.github.io/peu-sunass-economia/}

% Icono discreto de Instagram en TikZ
\newcommand{\igicon}[1][white!88]{%
  \tikz[baseline=-0.25ex]{%
    \draw[#1, line width=0.65pt, rounded corners=1.1pt] (0,0) rectangle (0.23cm,0.23cm);
    \draw[#1, line width=0.6pt] (0.115cm,0.115cm) circle (0.055cm);
    \fill[#1] (0.18cm,0.18cm) circle (0.014cm);
  }%
}

% Icono discreto de TikTok en TikZ (nota musical estilizada)
\newcommand{\tiktokicon}[1][white!88]{%
  \tikz[baseline=-0.25ex]{%
    \draw[#1, line width=0.82pt, line cap=round, line join=round]
      (0.11cm,0.06cm) -- (0.11cm,0.22cm)
      .. controls (0.12cm,0.16cm) and (0.16cm,0.145cm) .. (0.20cm,0.145cm);
    \draw[#1, line width=0.82pt, line cap=round]
      (0.11cm,0.065cm) arc[start angle=0, end angle=295, radius=0.048cm];
  }%
}

% Enlaces compactos de redes sociales: \backboneSocials{color}{tamaño}
\newcommand{\backboneSocials}[2]{%
  #2\color{#1}%
  \href{\backboneIgUrl}{\igicon[#1]}%
  \hspace{3.5pt}%
  \href{\backboneTtUrl}{\tiktokicon[#1]}%
  \hspace{4pt}%
  \href{\backboneIgUrl}{\textbf{@backbone.academy}}%
}


% Pie de página institucional en todas las páginas interiores
\lfoot{\scriptsize\color{azulOscuro}\textbf{BackBone Academy}\ \color{grisTexto}· \href{\backboneIgUrl}{\igicon[azulPrincipal]\ \textbf{@backbone.academy}} · \href{\backboneTtUrl}{\tiktokicon[azulPrincipal]\ \textbf{TikTok}} · \href{\backboneWebUrl}{Simulador Web}}
\cfoot{}
\rfoot{\small\thepage}
\renewcommand{\footrulewidth}{0.3pt}

% Encabezado superior para la portada de cada banco
\newcommand{\backboneportadaheader}{%
  \vspace*{0.1cm}
  \noindent\includegraphics[height=0.72cm]{../teoria/backbone_logo_azul.png}%
  \hfill%
  \raisebox{0.22cm}{\backboneSocials{azulPrincipal}{\small}}%
  \par\vspace{0.9cm}%
}

% Pie institucional para la portada de cada banco
\newcommand{\backboneportadafooter}{%
  {\small\color{azulOscuro}\textbf{BackBone Academy} \color{black}· Material de preparación académica.}\\[0.18cm]
  {\small\color{azulPrincipal}\href{\backboneIgUrl}{\igicon[azulPrincipal]\ \textbf{Instagram}} \quad·\quad \href{\backboneTtUrl}{\tiktokicon[azulPrincipal]\ \textbf{TikTok}} \textbf{(@backbone.academy)} \quad·\quad \href{\backboneWebUrl}{\textbf{Simulador Web Interactivo}}}\\[0.18cm]
  {\small Edición 2026}%
}


\begin{document}

% ============================================================
% PORTADA
% ============================================================
\begin{titlepage}
\thispagestyle{empty}
\centering
\backboneportadaheader

{\Large\bfseries BANCO DE PREGUNTAS}\\[0.35cm]
{\large EXAMEN DE ADMISIÓN}\\[0.15cm]
{\large PROGRAMA DE EXTENSIÓN UNIVERSITARIA (PEU)}\\[0.15cm]
{\large SUPERINTENDENCIA NACIONAL DE SERVICIOS DE SANEAMIENTO}\\[0.05cm]
{\large SUNASS}\\[1.0cm]

\rule{\linewidth}{0.9pt}\\[0.55cm]
{\LARGE\bfseries MÓDULO X: FINANZAS CORPORATIVAS Y EVALUACIÓN DE PROYECTOS}\\[0.35cm]
{\large Especialidad: Economía}\\[0.55cm]
\rule{\linewidth}{0.9pt}\\[1.0cm]

\begin{tabular}{ll}
\toprule
Número de preguntas: & 90 \\
Bloque A: & 30 preguntas \\
Bloque B: & 35 preguntas \\
Bloque C: & 25 preguntas \\
Duración sugerida: & 150 minutos \\
Nivel: & Finanzas intermedias--avanzadas \\
Formato: & Opción múltiple (A, B, C, D) \\
\bottomrule
\end{tabular}

\vfill
\backboneportadafooter
\end{titlepage}

% ============================================================
% INSTRUCCIONES
% ============================================================
\section*{Instrucciones generales}

\begin{enumerate}[label=\arabic*.,leftmargin=1.6em,itemsep=3pt]
  \item El presente banco contiene \textbf{90 preguntas de opción múltiple}. Cada pregunta tiene una sola respuesta correcta.
  \item No se descuenta por respuesta incorrecta; se recomienda responder todas las preguntas.
  \item El tiempo sugerido de resolución es de \textbf{150 minutos}. Use calculadora financiera o científica no programable.
  \item Las preguntas están organizadas en tres bloques: fundamentos teóricos, aplicaciones cuantitativas y casos regulatorios del sector saneamiento.
  \item El \textbf{solucionario completo} y la \textbf{tabla de claves rápidas} se encuentran al final. Se recomienda resolver el examen antes de consultarlos.
\end{enumerate}

\section*{Estructura del banco}

\begin{center}
\renewcommand{\arraystretch}{1.25}
\begin{tabular}{@{}p{1.4cm}>{\raggedright\arraybackslash}p{7.6cm}cc@{}}
\toprule
\textbf{Bloque} & \textbf{Contenido} & \textbf{Preguntas} & \textbf{Peso} \\
\midrule
A & Fundamentos teóricos: valor del dinero en el tiempo, tasas nominales y efectivas (TEA), ecuación de Fisher, anualidades, estados financieros, ratios de liquidez, solvencia, rentabilidad (DuPont), criterios VAN, TIR, TIRM, Payback y WACC. & 1--30 & 33.3\% \\
B & Aplicaciones cuantitativas: cálculo de anualidades y perpetuidades, análisis de estados financieros, flujo de caja libre (FCFF), VAN, TIR, tasa de Fisher, Costo Anual Equivalente (CAE), modelo CAPM, betas apalancadas y tablas de amortización. & 31--65 & 38.9\% \\
C & Casos con implicancias regulatorias: WACC regulatorio de SUNASS, Base de Activos Regulatorios (BAR), Valor Nuevo de Reemplazo (VNR), evaluación social en Invierte.pe, Tasa Social de Descuento (8\%), contratos APP en PTAR y sostenibilidad de EPS. & 66--90 & 27.8\% \\
\bottomrule
\end{tabular}
\end{center}

\section*{Competencias evaluadas}

\noindent El banco evalúa la capacidad de dominar la matemática financiera y el valor del dinero en el tiempo, interpretar con rigor balances y estados de resultados mediante ratios de liquidez, apalancamiento y rentabilidad, formular flujos de caja libres de proyectos, aplicar criterios de decisión de inversiones (VAN, TIR, TIRM, B/C y CAE), calcular el costo de capital propio mediante el CAPM y determinar el Costo Promedio Ponderado de Capital (WACC) regulatorio para remunerar inversiones en infraestructura de agua y saneamiento bajo el marco de SUNASS e Invierte.pe.

\section*{Notación principal}

{\small
\renewcommand{\arraystretch}{1.0}
\begin{center}
\begin{tabular}{@{}ll@{\hspace{1.0cm}}ll@{}}
\toprule
$VP, VF$ & Valor Presente y Valor Futuro & $TEA, TNA$ & Tasa Efectiva y Tasa Nominal Anual \\
$VAN / NPV$ & Valor Actual Neto & $TIR / IRR$ & Tasa Interna de Retorno \\
$WACC$ & Costo Promedio Ponderado de Capital & $K_e, K_d$ & Costo de capital propio y costo de la deuda \\
$ROA, ROE$ & Rentabilidad sobre activos y patrimonio & $FCLP / FCFF$ & Flujo de caja libre del proyecto \\
$CAE / EAC$ & Costo Anual Equivalente & $VNR$ & Valor Nuevo de Reemplazo \\
$BAR / RAB$ & Base de Activos Regulatorios & $TSD$ & Tasa Social de Descuento \\
\bottomrule
\end{tabular}
\end{center}
}

\newpage
\setcounter{qnum}{0}

\section*{Bloque A. Fundamentos teóricos y conceptuales}

\noindent\textit{Preguntas 1 a 30. Evalúan la comprensión analítica del valor del dinero en el tiempo, regímenes de interés, TEA y ecuación de Fisher, anualidades y perpetuidades, estados financieros, ratios de liquidez, solvencia y rentabilidad (DuPont), flujos de caja libres y criterios de evaluación de proyectos (VAN, TIR, TIRM, Payback, CAE y WACC).}

\q{El principio del 'Valor del Dinero en el Tiempo' establece fundamentalmente que una unidad monetaria hoy vale más que una unidad monetaria en el futuro debido a:}
\begin{opts}
    \item existencia exclusiva de intermediarios bancarios con encaje legal positivo.
    \item mandato normativo discrecional fijado anualmente por el Banco Central.
    \item preferencia por consumo presente, inflación esperada y costo de oportunidad del capital.
    \item volatilidad cambiaria esperada respecto a monedas extranjeras líquidas.
\end{opts}

\q{En el régimen de interés compuesto frente al de interés simple:}
\begin{opts}
    \item los intereses devengados se capitalizan y generan nuevos intereses en periodos sucesivos.
    \item la tasa de crecimiento del capital acumulado sigue una trayectoria estrictamente lineal.
    \item en interés compuesto no se cumple la propiedad de aditividad temporal del flujo de caja.
    \item los intereses generados en cada periodo se retiran sin integrarse al capital inicial.
\end{opts}

\q{La Tasa Efectiva Anual (TEA) se distingue de la Tasa Nominal Anual (TNA) en que la TEA:}
\begin{opts}
    \item es fijada por mandato administrativo anual del Ministerio de Economía y Finanzas.
    \item ignora la frecuencia de capitalización subperiódica de intereses dentro del año fiscal.
    \item mide el costo efectivo anual incorporando capitalización ($TEA = (1 + TNA/m)^m - 1$).
    \item es siempre estrictamente menor que la TNA cuando la capitalización es mensual o diaria.
\end{opts}

\q{La Ecuación de Fisher exacta que relaciona la tasa de interés nominal ($i$), la tasa de interés real ($r$) y la tasa de inflación esperada ($\pi$) se expresa formalmente como:}
\begin{opts}
    \item $1 + r = \frac{1 + i}{\pi}$
    \item $1 + i = (1 + r)(1 + \pi)$
    \item $1 + i = \frac{1 + r}{1 + \pi}$
    \item $r = (1 + i)(1 + \pi) - 1$
\end{opts}

\q{El Valor Presente ($VP$) de una corriente perpetua de flujos de caja constantes e iguales a $C$, que se cobran a partir del final del primer periodo con una tasa de descuento constante $r$, es exactamente:}
\begin{opts}
    \item $VP = \frac{C}{(1 + r)^n}$
    \item $VP = \frac{C}{r - g}$
    \item $VP = C \times r$
    \item $VP = \frac{C}{r}$
\end{opts}

\q{El modelo de Gordon y Shapiro para el Valor Presente de una perpetuidad creciente a una tasa constante $g$ anual a perpetuidad ($C_1$ a cobrar en $t=1$, con $r > g$) establece que:}
\begin{opts}
    \item $VP = C_1 \times (r - g)$
    \item $VP = \frac{C_1}{g}$
    \item $VP = \frac{C_1}{r - g}$
    \item $VP = \frac{C_1}{r + g}$
\end{opts}

\q{La relación matemática entre el Valor Presente de una anualidad anticipada ($VP_{ant}$) y el Valor Presente de una anualidad vencida u ordinaria ($VP_{venc}$) con idéntico número de cuotas y tasa de interés $r$ es:}
\begin{opts}
    \item $VP_{ant} = VP_{venc} \times (1 - r)$
    \item $VP_{ant} = VP_{venc} \times (1 + r)$
    \item $VP_{ant} = VP_{venc} \times r$
    \item $VP_{ant} = VP_{venc} / (1 + r)$
\end{opts}

\q{La Curva de Rendimiento (Yield Curve) que representa la estructura temporal de las tasas de interés suele tener pendiente positiva en condiciones macroeconómicas normales porque:}
\begin{opts}
    \item los inversionistas institucionales prefieren inmovilizar capital a largo plazo por liquidez.
    \item el Banco Central impone por norma una estructura temporal de tasas estrictamente creciente.
    \item los plazos largos conllevan mayor riesgo de tasa e inflación, exigiendo prima de liquidez.
    \item los emisores soberanos de largo plazo carecen de pasivos contingentes en su balance general.
\end{opts}

\q{En el Balance General de una empresa, la ecuación contable básica que expresa la igualdad fundamental de fuentes y usos de recursos es:}
\begin{opts}
    \item Activo = Pasivo + Patrimonio Neto
    \item Pasivo = Activo + Patrimonio Neto
    \item Patrimonio = Activo + Pasivo Neto
    \item Activo = Pasivo - Patrimonio Neto
\end{opts}

\q{La 'Prueba Ácida' (o ratio de liquidez severa) se diferencia de la 'Razón Corriente' en que la prueba ácida:}
\begin{opts}
    \item excluye los inventarios del activo corriente por ser el rubro realizable menos líquido.
    \item deduce los préstamos bancarios de largo plazo del pasivo corriente.
    \item multiplica las cuentas por cobrar comerciales por el margen bruto anual.
    \item suma los activos fijos depreciables al disponible neto de caja y bancos.
\end{opts}

\q{El 'Capital de Trabajo Neto' (Net Working Capital) de una empresa prestadora de servicios se define operativamente como:}
\begin{opts}
    \item el flujo de caja neto operativo generado durante el ejercicio contable.
    \item la diferencia entre Activo Corriente y Pasivo Corriente ($AC - PC$).
    \item la suma de las utilidades retenidas más las reservas legales del capital.
    \item el total de activos fijos tangibles netos de depreciación acumulada.
\end{opts}

\q{El ratio de Apalancamiento Financiero Deuda/Patrimonio ($D/E$) mide:}
\begin{opts}
    \item el financiamiento ajeno frente a fondos propios, reflejando el riesgo financiero.
    \item el número de veces que el capital de trabajo neto rota respecto a las ventas anuales totales.
    \item la rentabilidad operativa neta obtenida por cada sol invertido en activos intangibles.
    \item la proporción de activos corrientes que han sido liquidados en efectivo al cierre fiscal.
\end{opts}

\q{El ratio de Cobertura de Intereses (Times Interest Earned) se calcula como:}
\begin{opts}
    \item $\frac{EBITDA}{Activo\ Corriente}$
    \item $\frac{EBIT}{Gastos\ Financieros}$
    \item $\frac{Utilidad\ Neta}{Pasivo\ Total}$
    \item $\frac{Ventas}{Gastos\ Financieros}$
\end{opts}

\q{La diferencia conceptual entre la Rentabilidad sobre los Activos ($ROA$) y la Rentabilidad sobre el Patrimonio ($ROE$) radica en que:}
\begin{opts}
    \item el ROE mide exclusivamente el impacto de la inflación sectorial sobre el capital de trabajo de la empresa.
    \item el ROA se aplica únicamente a bancos comerciales mientras el ROE mide liquidez en empresas reales.
    \item el ROA supera siempre al ROE cuando la firma contrata deuda bancaria con escudo fiscal positivo.
    \item el ROA evalúa la eficiencia de los activos totales y el ROE el retorno sobre el patrimonio propio.
\end{opts}

\q{La descomposición DuPont del $ROE$ desagrega la rentabilidad sobre el patrimonio en el producto de tres motores financieros clave:}
\begin{opts}
    \item Margen bruto operativo $\times$ Razón corriente $\times$ Inflación esperada anual.
    \item Margen Neto $\times$ Rotación de Activos $\times$ Multiplicador de Capital ($A/P$).
    \item Costo de ventas $\times$ Tasa impositiva efectiva $\times$ Cobertura de deuda ($EBIT$).
    \item Flujo de caja libre $\times$ Tasa de interés nominal $\times$ Costo promedio de capital.
\end{opts}

\q{El 'Ciclo de Conversión del Efectivo' (Cash Conversion Cycle) mide el tiempo transcurrido desde que la empresa invierte efectivo en insumos hasta que recupera el efectivo por cobranza, calculándose como:}
\begin{opts}
    \item Días de inventario + Días de cobro - Días de pago.
    \item Rotación de inventarios + Días de pago / 365.
    \item Días de cobro - Días de inventario + Días de pago.
    \item Días de cobro + Periodo promedio de pago a proveedores.
\end{opts}

\q{Al proyectar el Flujo de Caja Libre del Proyecto (FCLP o Free Cash Flow to Firm - FCFF):}
\begin{opts}
    \item suma la utilidad neta contable después de intereses más el total de pasivos no corrientes.
    \item descuenta íntegramente las cuotas de amortización de préstamos y los gastos financieros.
    \item excluye el servicio de la deuda (intereses y amortización), midiendo el valor operativo puro.
    \item incluye como egresos relevantes los costos hundidos devengados en ejercicios anteriores.
\end{opts}

\q{En la preparación del flujo de caja de un proyecto de inversión, un 'Costo Hundido' (Sunk Cost), como el dinero ya gastado en un estudio geológico previo no reembolsable:}
\begin{opts}
    \item debe excluirse del flujo incremental por ser un gasto histórico irrecuperable.
    \item debe registrarse como inversión inicial negativa en el periodo $t=0$.
    \item debe actualizarse al costo de oportunidad e incluirse como flujo operativo.
    \item debe depreciarse tributariamente de forma lineal durante diez años.
\end{opts}

\q{El criterio del Valor Actual Neto (VAN / NPV) para la toma de decisiones de inversión establece que un proyecto financieramente independiente debe aceptarse si:}
\begin{opts}
    \item la relación Beneficio/Costo resulta estrictamente menor que la unidad ($B/C < 1$).
    \item $VAN < 0$, indicando una menor salida de capital respecto al costo ponderado.
    \item el valor actual de los flujos iguala al costo inicial ($VAN = 0$) en sector privado.
    \item $VAN > 0$, pues los flujos descontados superan la inversión y crean riqueza.
\end{opts}

\q{Una propiedad matemática y financiera fundamental que hace que el criterio del VAN sea superior a otros indicadores es la 'Aditividad de Valor', la cual significa que:}
\begin{opts}
    \item el VAN de todo proyecto aumenta monótonamente al subir el costo de capital.
    \item para proyectos independientes, se cumple $VAN(A+B) = VAN(A) + VAN(B)$.
    \item la TIR conjunta de un portafolio equivale a la suma de las TIR individuales.
    \item el VAN resulta insensible a los flujos de caja de horizontes lejanos.
\end{opts}

\q{La Tasa Interna de Retorno (TIR / IRR) se define rigurosamente como:}
\begin{opts}
    \item la tasa de descuento que anula el Valor Actual Neto ($VAN(TIR) = 0$).
    \item la tasa de interés pasiva fijada por el banco comercial más grande del mercado.
    \item el cociente contable promedio entre la utilidad bruta y los activos totales.
    \item la tasa que maximiza la suma aritmética simple de flujos de caja nominales.
\end{opts}

\q{Una limitación técnica crítica de la Tasa Interna de Retorno frente al VAN radica en el supuesto implícito de reinversión de los flujos intermedios, según el cual:}
\begin{opts}
    \item la TIR asume reinversión a tasa cero mientras el VAN supone ausencia total de reinversión.
    \item el VAN asume reinversión al costo de capital ($WACC$), mientras la TIR asume reinversión a la propia TIR.
    \item ambos métodos suponen reinversión obligatoria a la tasa libre de riesgo de los bonos soberanos.
    \item el VAN asume reinversión a la tasa de inflación; la TIR asume reinversión al costo de oportunidad del capital.
\end{opts}

\q{En proyectos con flujos de caja no convencionales (donde ocurren múltiples cambios de signo en los flujos a lo largo del tiempo, por ejemplo por costos de cierre de una mina o planta de tratamiento), la Regla de los Signos de Descartes implica que:}
\begin{opts}
    \item el proyecto carece de Valor Actual Neto definido para tasas reales positivas.
    \item el proyecto debe rechazarse automáticamente sin calcular su Valor Actual Neto.
    \item el VAN es estrictamente cóncavo y creciente respecto a la tasa de descuento relevante.
    \item pueden existir múltiples TIR (tantas como cambios de signo), invalidando ese criterio.
\end{opts}

\q{Al comparar dos proyectos de inversión mutuamente excluyentes con diferente escala o diferente perfil temporal de flujos (donde las curvas de VAN se intersectan en la llamada Tasa de Fisher), el criterio que debe prevalecer siempre es:}
\begin{opts}
    \item el proyecto que requiera la mayor inversión inicial en infraestructura física.
    \item el Periodo de Recuperación simple más corto para minimizar el riesgo operativo.
    \item la Tasa Interna de Retorno (TIR), eligiendo el proyecto con mayor porcentaje de rentabilidad.
    \item el Valor Actual Neto (VAN) al costo de capital, por maximizar el valor absoluto.
\end{opts}

\q{La Tasa Interna de Retorno Modificada (TIRM / MIRR) supera las limitaciones de la TIR tradicional porque:}
\begin{opts}
    \item prescinde de conocer la vida útil del proyecto y la estructura de financiamiento de la firma.
    \item actualiza egresos al financiamiento y capitaliza ingresos al costo de capital, dando tasa única.
    \item elimina del flujo los desembolsos negativos iniciales para asegurar convergencia numérica.
    \item transforma automáticamente cualquier secuencia de flujos para que arroje un VAN positivo.
\end{opts}

\q{El método del Periodo de Recuperación del Capital (Payback simple) presenta dos deficiencias teóricas fundamentales para la toma de decisiones óptimas:}
\begin{opts}
    \item sobrestima sistemáticamente el capital de trabajo y la depreciación.
    \item ignora el valor del dinero en el tiempo y los flujos posteriores al recupero.
    \item exige cálculo matricial iterativo y depende de la deuda contable.
    \item asume la liquidación forzosa de activos fijos al cumplir el plazo.
\end{opts}

\q{El Costo Anual Equivalente (CAE / EAC) es la herramienta financiera idónea para:}
\begin{opts}
    \item comparar proyectos mutuamente excluyentes con distintas vidas útiles y costos.
    \item medir el impacto de la inflación acumulada sobre los contratos de construcción civil.
    \item evaluar la solvencia de corto plazo de una firma ante choques de capital de trabajo.
    \item determinar el precio de emisión de acciones preferentes en el mercado secundario.
\end{opts}

\q{La fórmula clásica del Costo Promedio Ponderado de Capital (WACC / CPPC) calculada sobre la estructura de capital a valores de mercado es:}
\begin{opts}
    \item $WACC = K_e (1 - T) \left(\frac{E}{V}\right) + K_d \left(\frac{D}{V}\right)$
    \item $WACC = K_e \left(\frac{E}{V}\right) + K_d (1 - T) \left(\frac{D}{V}\right)$
    \item $WACC = K_e \left(\frac{E}{V}\right) + \frac{K_d}{(1 - T)} \left(\frac{D}{V}\right)$
    \item $WACC = K_e \left(\frac{E}{V}\right) + K_d \left(\frac{D}{V}\right)$
\end{opts}

\q{En el modelo Capital Asset Pricing Model (CAPM) para el cálculo del costo del capital propio ($K_e$), la ecuación representativa es $K_e = R_f + \beta (R_m - R_f)$. El parámetro Beta ($\beta$) mide:}
\begin{opts}
    \item el riesgo diversificable o específico propio de la gestión interna de la firma.
    \item la volatilidad idiosincrásica del tipo de cambio frente al oro monetario.
    \item el riesgo sistemático no diversificable respecto al portafolio de mercado.
    \item la rentabilidad nominal libre de riesgo de los bonos soberanos del tesoro.
\end{opts}

\q{La Proposición I de Modigliani y Miller (1958) en un mundo ideal de mercados de capitales perfectos (sin impuestos corporativos, sin costos de quiebra y con información simétrica) establece que:}
\begin{opts}
    \item las empresas deben financiarse exclusivamente con utilidades retenidas.
    \item el valor de mercado crece linealmente al aumentar el ratio de deuda.
    \item el valor de mercado de la firma es totalmente independiente de su estructura de capital.
    \item el costo promedio ponderado de capital equivale a la tasa libre de riesgo.
\end{opts}


\section*{Bloque B. Aplicaciones cuantitativas y financieras}

\noindent\textit{Preguntas 31 a 65. Evalúan la capacidad matemática para calcular tasas efectivas, anualidades, ratios financieros, flujos de caja libres del proyecto (FCLP), VAN, TIR, tasa de Fisher, Costo Anual Equivalente (CAE), WACC corporativo, modelo CAPM y cronogramas de amortización de deuda.}

\q{Un banco comercial aprueba un préstamo a una EPS a una Tasa Nominal Anual del $TNA = 12\%$ con capitalización mensual de intereses. ¿Cuál es la Tasa Efectiva Anual (TEA) equivalente que realmente pagará la empresa?}
\begin{opts}
    \item $TEA = 13.55\%$
    \item $TEA = 11.45\%$
    \item $TEA = 12.00\%$
    \item $TEA \approx 12.68\%$
\end{opts}

\q{Si un depósito a plazo fijo en una entidad financiera rinde una tasa de interés nominal del $i = 8.15\%$ anual y la inflación proyectada para el año es del $\pi = 3.0\%$, la tasa de interés real ($r$) exacta según la ecuación de Fisher es:}
\begin{opts}
    \item $r = 5.15\%$
    \item $r = 11.39\%$
    \item $r = 2.44\%$
    \item $r = 5.00\%$
\end{opts}

\q{Una Planta de Tratamiento genera ahorros netos de energía constantes de 50000 soles anuales que se proyectan a perpetuidad. Si la tasa de descuento corporativa relevante es del $r = 8.0\%$ anual, ¿cuál es el Valor Presente ($VP$) de dichos ahorros?}
\begin{opts}
    \item $VP = 625000$ soles
    \item $VP = 500000$ soles
    \item $VP = 400000$ soles
    \item $VP = 800000$ soles
\end{opts}

\q{Un proyecto hidroeléctrico que abastece a una EPS proyecta un flujo de caja de $C_1 = 30000$ soles para el primer año, el cual crecerá a una tasa anual constante del $g = 3.0\%$ indefinidamente. Si la tasa de descuento es $r = 9.0\%$, el valor presente de la perpetuidad creciente es:}
\begin{opts}
    \item $VP = 1000000$ soles
    \item $VP = 500000$ soles
    \item $VP = 333333$ soles
    \item $VP = 600000$ soles
\end{opts}

\q{Una EPS adquiere un equipo de laboratorio mediante un préstamo de 100000 soles a amortizar en 5 cuotas anuales fijas vencidas iguales a una tasa de interés del $10\%$ anual. El factor de recuperación del capital es $\left[\frac{0.10(1.10)^5}{(1.10)^5 - 1}\right] \approx 0.263797$. La cuota fija anual ($C$) a pagar es:}
\begin{opts}
    \item $C = 24500.50$ soles
    \item $C = 26379.75$ soles
    \item $C = 20000.00$ soles
    \item $C = 30000.00$ soles
\end{opts}

\q{El Balance General de una empresa prestadora registra: Activo Corriente $= 160000$ soles (de los cuales $60000$ soles son Inventarios de repuestos) y Pasivo Corriente $= 80000$ soles. Los ratios de Razón Corriente y Prueba Ácida son respectivamente:}
\begin{opts}
    \item Razón Corriente $= 1.5$ veces; Prueba Ácida $= 1.0$ veces
    \item Razón Corriente $= 2.5$ veces; Prueba Ácida $= 1.75$ veces
    \item Razón Corriente $= 2.0$ veces; Prueba Ácida $= 1.25$ veces
    \item Razón Corriente $= 1.8$ veces; Prueba Ácida $= 0.8$ veces
\end{opts}

\q{Una EPS reporta Activos Totales por 500000 soles y Pasivos Totales exigibles por 300000 soles. ¿Cuál es su Patrimonio Neto y su ratio de Apalancamiento Deuda/Patrimonio ($D/E$)?}
\begin{opts}
    \item Patrimonio $= 200000$ soles y $D/E = 1.50$ (150\%)
    \item Patrimonio $= 150000$ soles y $D/E = 2.00$ (200\%)
    \item Patrimonio $= 200000$ soles y $D/E = 0.60$ (60\%)
    \item Patrimonio $= 300000$ soles y $D/E = 1.00$ (100\%)
\end{opts}

\q{Una empresa reporta Utilidad Operativa antes de intereses e impuestos ($EBIT$) de 120000 soles y registra gastos por intereses de su deuda financiera de 30000 soles en el año. Su ratio de Cobertura de Intereses es:}
\begin{opts}
    \item Cobertura $= 4.00$ veces
    \item Cobertura $= 0.25$ veces
    \item Cobertura $= 5.00$ veces
    \item Cobertura $= 3.00$ veces
\end{opts}

\q{Mediante el sistema DuPont, una empresa presenta: Margen de Utilidad Neta $= 8.0\%$, Rotación de Activos Totales $= 1.5$ veces y Multiplicador del Capital (Activo/Patrimonio) $= 1.50$. ¿Cuál es el $ROE$ de la empresa?}
\begin{opts}
    \item $ROE = 15.0\%$
    \item $ROE = 9.5\%$
    \item $ROE = 18.0\%$
    \item $ROE = 12.0\%$
\end{opts}

\q{Una EPS registra en su gestión operativa: Periodo Promedio de Cobro a usuarios $= 45$ días, Días de Rotación de Inventarios en almacén $= 60$ días y Periodo Promedio de Pago a sus proveedores químicos $= 40$ días. ¿Cuál es su Ciclo de Conversión del Efectivo ($CCE$)?}
\begin{opts}
    \item $CCE = 75$ días
    \item $CCE = 65$ días
    \item $CCE = 55$ días
    \item $CCE = 145$ días
\end{opts}

\q{En la evaluación de un proyecto de microcuencas, la Utilidad Operativa proyectada para el año 1 es $EBIT = 200000$ soles. La tasa de impuesto a la renta es $T = 30\%$, la depreciación del periodo es $40000$ soles, la inversión en activos fijos es $CAPEX = 50000$ soles y la variación en capital de trabajo neto es $\Delta CTN = 10000$ soles. El Flujo de Caja Libre del Proyecto ($FCLP$) para el año 1 es:}
\begin{opts}
    \item $FCLP = 140000$ soles
    \item $FCLP = 180000$ soles
    \item $FCLP = 120000$ soles
    \item $FCLP = 100000$ soles
\end{opts}

\q{Un proyecto de sectorización de redes exige una inversión inicial de $I_0 = 100000$ soles en $t=0$ y genera flujos de caja netos de $40000$ soles al final de cada uno de los siguientes 3 años ($t=1, 2, 3$). La tasa de descuento del costo de capital es $r = 10\%$ anual (factor de anualidad $A_3^{10\%} \approx 2.48685$). El Valor Actual Neto ($VAN$) del proyecto es:}
\begin{opts}
    \item $VAN = +20000$ soles
    \item $VAN = -526$ soles
    \item $VAN = 0$ soles
    \item $VAN = +526$ soles
\end{opts}

\q{Un proyecto pequeño requiere una inversión inicial de $I_0 = 10000$ soles y genera flujos de $6000$ soles en el año 1 y $6000$ soles en el año 2. La Tasa Interna de Retorno ($TIR$) del proyecto satisface $10000 = \frac{6000}{1+TIR} + \frac{6000}{(1+TIR)^2}$. El valor aproximado de la TIR es:}
\begin{opts}
    \item $TIR \approx 20.00\%$
    \item $TIR \approx 13.07\%$
    \item $TIR \approx 10.00\%$
    \item $TIR \approx 16.67\%$
\end{opts}

\q{Considere dos proyectos de inversión mutuamente excluyentes: Proyecto A ($I_0 = -1000$, flujos: $t_1 = 700, t_2 = 700$) y Proyecto B ($I_0 = -1000$, flujos: $t_1 = 200, t_2 = 1300$). La 'Tasa de Fisher' que iguala el VAN de ambos proyectos se obtiene resolviendo $VAN(A) - VAN(B) = 0$. Esta tasa de corte es:}
\begin{opts}
    \item $r_{Fisher} = 10.0\%$
    \item $r_{Fisher} = 8.0\%$
    \item $r_{Fisher} = 20.0\%$
    \item $r_{Fisher} = 15.0\%$
\end{opts}

\q{Un sistema de bombeo solar cuesta 70000 soles y genera ahorros de combustible de 25000 soles anuales constantes durante 5 años. El Periodo de Recuperación simple ($Payback$) del capital es:}
\begin{opts}
    \item $Payback = 4.00$ años
    \item $Payback = 3.50$ años
    \item $Payback = 2.80$ años
    \item $Payback = 2.00$ años
\end{opts}

\q{En la evaluación de un proyecto de control de pérdidas de agua, el Valor Presente de los beneficios económicos totales asciende a $VP(B) = 280000$ soles y el Valor Presente de los costos totales asciende a $VP(C) = 200000$ soles. La Relación Beneficio/Costo ($B/C$) y el $VAN$ son:}
\begin{opts}
    \item $B/C = 1.20$ y $VAN = +40000$ soles
    \item $B/C = 0.71$ y $VAN = -80000$ soles
    \item $B/C = 1.50$ y $VAN = +100000$ soles
    \item $B/C = 1.40$ y $VAN = +80000$ soles
\end{opts}

\q{Para renovar un sistema de telemetría de reservorios se evalúan dos bombas: Bomba Alfa tiene un Valor Presente de Costos de $50000$ soles y vida útil de 3 años ($A_3^{10\%} = 2.48685$). Bomba Beta tiene un Valor Presente de Costos de $75000$ soles y vida útil de 5 años ($A_5^{10\%} = 3.79079$). Con una tasa de descuento del $10\%$, compare sus Costos Anuales Equivalentes ($CAE$):}
\begin{opts}
    \item $CAE_{Alfa} = 15000.0$ soles vs $CAE_{Beta} = 18000.0$ soles
    \item $CAE_{Alfa} = 20105.7$ soles vs $CAE_{Beta} = 19784.8$ soles
    \item $CAE_{Alfa} = 25000.0$ soles vs $CAE_{Beta} = 22000.0$ soles
    \item $CAE_{Alfa} = 21500.0$ soles vs $CAE_{Beta} = 21500.0$ soles
\end{opts}

\q{Una empresa de servicios se financia con $60\%$ de capital propio ($E/V = 0.60$) y $40\%$ de deuda bancaria ($D/V = 0.40$). El costo del capital propio es $K_e = 12.0\%$, la tasa de interés de la deuda es $K_d = 7.0\%$ y la tasa del impuesto a la renta es $T = 30\%$. ¿Cuál es el Costo Promedio Ponderado de Capital ($WACC$)?}
\begin{opts}
    \item $WACC = 9.16\%$
    \item $WACC = 11.20\%$
    \item $WACC = 8.50\%$
    \item $WACC = 10.00\%$
\end{opts}

\q{Aplicando el modelo CAPM adaptado a mercados emergentes ($K_e = R_f + \beta (R_m - R_f) + RP$): la tasa libre de riesgo es $R_f = 4.0\%$, la prima de riesgo de mercado es $(R_m - R_f) = 6.0\%$, la Beta apalancada del sector es $\beta = 0.75$ y el spread de Riesgo País (EMBI+) es $RP = 1.50\%$. El costo del capital propio ($K_e$) es:}
\begin{opts}
    \item $K_e = 11.50\%$
    \item $K_e = 9.25\%$
    \item $K_e = 10.00\%$
    \item $K_e = 8.50\%$
\end{opts}

\q{Una corporación mantiene una deuda financiera perpetua de $D = 3000000$ soles a una tasa de interés del $6.0\%$. Si la tasa del impuesto a la renta corporativa es $T = 30\%$, ¿cuál es el Valor Presente del Escudo Fiscal ($VP_{EF}$) según Modigliani y Miller?}
\begin{opts}
    \item $VP_{EF} = 90000$ soles
    \item $VP_{EF} = 900000$ soles
    \item $VP_{EF} = 180000$ soles
    \item $VP_{EF} = 54000$ soles
\end{opts}

\q{En el sistema de amortización francés de una deuda de 100000 soles a 4 años con cuotas anuales fijas de 31547.08 soles al $10\%$ de interés anual, ¿cuánto corresponde a intereses y cuánto a amortización del principal en la primera cuota?}
\begin{opts}
    \item Interés $= 10000.00$ soles y Amortización $= 21547.08$ soles
    \item Interés $= 3154.71$ soles y Amortización $= 28392.37$ soles
    \item Interés $= 5000.00$ soles y Amortización $= 26547.08$ soles
    \item Interés $= 15000.00$ soles y Amortización $= 16547.08$ soles
\end{opts}

\q{En el sistema de amortización constante alemán de una deuda de 60000 soles a pagar en 3 cuotas anuales al $10\%$ de interés anual sobre saldos:}
\begin{opts}
    \item todas las cuotas totales a pagar resultan idénticas durante los tres años.
    \item no se devengan intereses financieros hasta la cancelación de la última cuota.
    \item la amortización del principal es creciente mientras el interés permanece fijo.
    \item la amortización del principal es fija ($20000$ soles) y las cuotas decrecen.
\end{opts}

\q{Un bono corporativo cupón cero con valor nominal de 1000 soles vence dentro de exactamente 2 años. Si la tasa de rendimiento requerida al vencimiento ($YTM$) es del $8\%$ anual, el precio de mercado actual del bono es:}
\begin{opts}
    \item $P = 920.00$ soles
    \item $P = 857.34$ soles
    \item $P = 840.00$ soles
    \item $P = 1000.00$ soles
\end{opts}

\q{Si un bono tiene una Duración Modificada de $D_{mod} = 4.5$ años y la tasa de interés de mercado sube en 50 puntos básicos ($+0.50\%$), la variación porcentual aproximada en el precio del bono es:}
\begin{opts}
    \item $+0.50\%$
    \item $+2.25\%$
    \item $-2.25\%$
    \item $-4.50\%$
\end{opts}

\q{Un proyecto presenta una inversión de 1000 soles y genera flujos de $600$ soles en año 1 y $800$ soles en año 2. Con una tasa de reinversión del $10\%$, el Valor Futuro terminal en el año 2 es $VF = 600(1.10) + 800 = 660 + 800 = 1460$ soles. La Tasa Interna de Retorno Modificada ($TIRM$) a 2 años es:}
\begin{opts}
    \item $TIRM = 10.00\%$
    \item $TIRM \approx 20.83\%$
    \item $TIRM = 23.00\%$
    \item $TIRM = 15.50\%$
\end{opts}

\q{Una empresa tiene costos fijos operativos de 120000 soles al año. Vende cada metro cúbico tratado a $P = 5.0$ soles y su costo variable unitario es $CV = 2.0$ soles. El punto de equilibrio operativo ($Q_{equilibrio}$) en unidades físicas es:}
\begin{opts}
    \item $Q = 60000 \; m^3$
    \item $Q = 40000 \; m^3$
    \item $Q = 24000 \; m^3$
    \item $Q = 30000 \; m^3$
\end{opts}

\q{El Grado de Apalancamiento Operativo (GAO) de una empresa es $GAO = 2.5$. Si las ventas de la compañía aumentan en un $8.0\%$, ¿en qué porcentaje se incrementará su Utilidad Operativa ($EBIT$)?}
\begin{opts}
    \item $10.5\%$
    \item $16.0\%$
    \item $8.0\%$
    \item $20.0\%$
\end{opts}

\q{Si una firma tiene un Grado de Apalancamiento Operativo de $GAO = 2.0$ y un Grado de Apalancamiento Financiero de $GAF = 1.5$, su Grado de Apalancamiento Total ($GAT$) es:}
\begin{opts}
    \item $GAT = 3.50$
    \item $GAT = 1.33$
    \item $GAT = 0.50$
    \item $GAT = 3.00$
\end{opts}

\q{En el método del Valor Presente Ajustado (APV), un proyecto financiado en parte con deuda subsidiada por el Estado se evalúa sumando:}
\begin{opts}
    \item la TIR ajustada del proyecto multiplicada por el ratio de deuda contable de la firma.
    \item el costo histórico de reposición menos el flujo de caja libre operativo del periodo.
    \item el $VAN_{base}$ sin deuda más el valor presente de escudos fiscales y subsidios.
    \item el beneficio neto contable ajustado por la depreciación y el EBIT sin impuestos.
\end{opts}

\q{Bajo el enfoque de opciones reales, una EPS que diseña una planta de agua con espacio físico disponible para duplicar su capacidad en el futuro si la demanda crece está adquiriendo:}
\begin{opts}
    \item un pasivo contingente de largo plazo que reduce el valor presente neto base.
    \item un contrato forward de cobertura sobre precios de tuberías de PVC y bombeo.
    \item una opción real de abandono financiero equivalente a una opción Put europea.
    \item una Opción Real de Expansión (tipo Call) que añade valor económico al VAN.
\end{opts}

\q{Si el costo del capital desapalancado de una industria es $K_u = 10.0\%$, el costo de la deuda es $K_d = 6.0\%$, la tasa impositiva es $T = 30\%$ y la relación deuda/patrimonio es $D/E = 0.50$, según la Proposición II de Modigliani-Miller con impuestos, el costo de capital propio apalancado $K_e$ es:}
\begin{opts}
    \item $K_e = 12.00\%$
    \item $K_e = 8.60\%$
    \item $K_e = 10.00\%$
    \item $K_e = 11.40\%$
\end{opts}

\q{Al proyectar flujos de caja en presencia de inflación, la regla de consistencia financiera básica exige que:}
\begin{opts}
    \item se excluya por completo la depreciación fiscal de activos de los flujos de caja nominales.
    \item la tasa de descuento nominal equivalga exactamente al doble de la tasa de inflación anual.
    \item los flujos nominales se descuenten a tasa nominal y los reales a tasa real, dando igual VAN.
    \item los flujos de caja nominales se descuenten siempre con la tasa de descuento real sin ajuste.
\end{opts}

\q{Al término del horizonte de evaluación de 10 años de una planta de tratamiento, el Valor Residual (o valor terminal) estimado como perpetuidad del flujo normalizado del año 10 ($FCLP_{10} = 80000$ soles) descontado a una tasa de $r = 10\%$ sin crecimiento ($g=0$) es:}
\begin{opts}
    \item $VR_{10} = 0$ soles al culminar la vida útil en $t=10$
    \item $VR_{10} = 80000$ soles en $t=10$, descontado a $t=0$
    \item $VR_{10} = 400000$ soles en $t=10$, descontado a $t=0$
    \item $VR_{10} = 800000$ soles en $t=10$, descontado a $t=0$
\end{opts}

\q{Si una acción de una empresa de servicios cotiza a 30 soles y reporta una utilidad neta por acción de 2.50 soles, su ratio Precio/Beneficio (P/E) es:}
\begin{opts}
    \item $P/E = 12.0$ veces
    \item $P/E = 8.33$ veces
    \item $P/E = 5.0$ veces
    \item $P/E = 75.0$ veces
\end{opts}

\q{Un bono cero cupón a 3 años de valor nominal 1000 soles se cotiza hoy en el mercado a 751.31 soles. ¿Cuál es su Rendimiento al Vencimiento ($YTM$)?}
\begin{opts}
    \item $YTM = 10.00\%$
    \item $YTM = 12.00\%$
    \item $YTM = 8.00\%$
    \item $YTM = 6.50\%$
\end{opts}


\section*{Bloque C. Casos con implicancias regulatorias y sector saneamiento}

\noindent\textit{Preguntas 66 a 90. Evalúan aplicaciones al marco tarifario y financiero de SUNASS y saneamiento peruano: WACC regulatorio, Base de Activos Regulatorios (BAR), Valor Nuevo de Reemplazo (VNR), evaluación social en Invierte.pe (Tasa Social de Descuento del 8\%), contratos APP en PTAR, descalce cambiario y sostenibilidad de EPS.}

\q{En la metodología regulatoria de SUNASS para la determinación del Costo Promedio Ponderado de Capital (WACC regulatorio) aplicado a las tarifas quinquenales de agua:}
\begin{opts}
  \item se asume que el costo del capital propio es igual a la inflación del INEI.
  \item se fija una estructura de capital regulatoria de referencia normativa (típicamente 50\% deuda y 50\% capital propio), incentivando a la EPS a optimizar su financiamiento sin trasladar ineficiencias de endeudamiento a los usuarios.
  \item se utiliza la estructura contable de deuda histórica de cada municipio.
  \item se prohíbe el uso de endeudamiento bancario.
\end{opts}

\q{Para estimar la Tasa Libre de Riesgo ($R_f$) en el CAPM del WACC de las EPS, SUNASS toma como parámetro de referencia internacional:}
\begin{opts}
  \item la tasa de inflación de Lima Metropolitana.
  \item la tasa de las cuentas de ahorro del Banco de la Nación.
  \item el rendimiento promedio de los Bonos del Tesoro de los Estados Unidos (US Treasury Bonds) a 10 o 30 años, por ser el activo con menor riesgo de crédito a nivel global.
  \item el rendimiento de los bonos corporativos de empresas mineras.
\end{opts}

\q{El spread de Riesgo País incorporado por SUNASS en el costo de oportunidad del capital propio de las EPS se mide usualmente a través de:}
\begin{opts}
  \item la tasa de desempleo juvenil de Lima.
  \item el arancel promedio a la importación de bienes de capital.
  \item el índice de precios al consumidor de Sedapal.
  \item el indicador EMBI+ Perú (Emerging Markets Bond Index) elaborado por J.P. Morgan, que refleja el diferencial de rendimiento entre los bonos soberanos peruanos y los bonos del tesoro norteamericano.
\end{opts}

\q{La Base de Activos Regulatorios (BAR / RAB) sobre la cual SUNASS aplica la tasa WACC para remunerar el capital en las tarifas de las EPS comprende:}
\begin{opts}
  \item el valor económico de los activos fijos netos en operación directamente afectos y utilizados en la prestación eficiente de los servicios de agua potable y alcantarillado.
  \item el total de activos intangibles y patentes de software de la empresa.
  \item el monto total de las cuentas por cobrar incobrables de usuarios morosos.
  \item las inversiones financieras especulativas en fondos mutuos.
\end{opts}

\q{Para valorar la infraestructura de redes subterráneas y plantas antiguas de agua, SUNASS y la teoría regulatoria aplican la metodología del Valor Nuevo de Reemplazo (VNR), la cual valora los activos según:}
\begin{opts}
  \item el precio de venta de las acciones en bolsa.
  \item el costo corriente que demandaría construir y poner en operación hoy un activo moderno equivalente con tecnología eficiente y capacidad adecuada, deduciendo la depreciación acumulada por estado de conservación.
  \item el valor de remate de la chatarra metálica.
  \item su costo histórico de adquisición contable registrado hace 40 años.
\end{opts}

\q{Aunque las 50 EPS del Perú son empresas públicas municipales que no buscan maximizar beneficios como una corporación privada, el regulador SUNASS les reconoce obligatoriamente una tasa WACC positiva en las tarifas para:}
\begin{opts}
  \item trasladar utilidades al Tesoro Público central.
  \item permitir el pago de bonos de productividad ilimitados a los funcionarios municipales.
  \item asegurar la sostenibilidad financiera a largo plazo, la reposición oportuna de activos depreciados y la capacidad de apalancar crédito para financiar la expansión de cobertura.
  \item financiar campañas de publicidad electoral en las provincias.
\end{opts}

\q{En el Sistema Nacional de Programación Multianual y Gestión de Inversiones (Invierte.pe), la evaluación de proyectos públicos de saneamiento se realiza bajo el enfoque de:}
\begin{opts}
  \item análisis de opciones financieras de derivados OTC.
  \item evaluación financiera privada exclusiva con precios de mercado.
  \item aprobación automática sin formulación de flujos de caja.
  \item evaluación social de proyectos, corrigiendo los precios de mercado mediante Factores de Corrección Social para reflejar el verdadero costo de oportunidad para la sociedad peruana.
\end{opts}

\q{La Tasa Social de Descuento (TSD) fijada por el Ministerio de Economía y Finanzas (MEF) en los anexos metodológicos de Invierte.pe para la evaluación de proyectos de saneamiento es actualmente de:}
\begin{opts}
  \item $8.0\%$ anual en términos reales
  \item $15.0\%$ anual
  \item $3.0\%$ anual
  \item $12.0\%$ anual
\end{opts}

\q{En la evaluación social de un proyecto de zanjas y reservorios en una zona andina deprimida, el Factor de Corrección Social de la Mano de Obra No Calificada (FC = 0.41 a 0.60 según región) es menor a 1 porque:}
\begin{opts}
  \item el Estado exonera del pago de remuneraciones a los campesinos.
  \item el salario de mercado supera el verdadero costo de oportunidad social del trabajador debido al desempleo o subempleo rural, por lo que el costo social para el país es menor al gasto contable de la planilla.
  \item los trabajadores rurales son más productivos que los ingenieros.
  \item la mano de obra se financia con donaciones externas.
\end{opts}

\q{Cuando un proyecto de agua potable rural en una comunidad nativa remota genera beneficios de reducción de parasitosis difíciles de cuantificar en soles, la metodología de evaluación de Invierte.pe que debe aplicarse es:}
\begin{opts}
  \item la TIR de mercado del proyecto.
  \item el Análisis Costo-Beneficio monetario con VAN social.
  \item el Análisis Costo-Eficacia o Costo-Efectividad, seleccionando la alternativa técnica que logre la meta de provisión y calidad al menor Costo Promedio por beneficiario o por metro cúbico.
  \item el remate público de la infraestructura.
\end{opts}

\q{Los bancos multilaterales (BID, Banco Mundial, KfW) que otorgan créditos de inversión a las EPS exigen el cumplimiento de un 'Ratio de Cobertura del Servicio de la Deuda' (DSCR), el cual se define como:}
\begin{opts}
  \item $DSCR = \frac{Pasivo\ Total}{Ventas\ Anuales}$
  \item $DSCR = \frac{Activo\ Fijo}{Patrimonio}$
  \item $DSCR = \frac{Utilidad\ Neta}{EBITDA}$
  \item $DSCR = \frac{Flujo\ de\ Caja\ Operativo\ Neto}{Servicio\ de\ la\ Deuda\ (Principal + Intereses)}$, exigiendo usualmente que sea $\ge 1.20$ veces para garantizar el repago.
\end{opts}

\q{El 'Periodo Promedio de Cobranza' de una EPS que supera los 120 días comerciales refleja un grave problema financiero ocasionado frecuentemente por:}
\begin{opts}
  \item la elevada morosidad y retrasos de pago sistemáticos de instituciones públicas (ministerios, comisarías, colegios estatales y hospitales) e industrias, inmovilizando capital de trabajo y forzando iliquidez.
  \item que los usuarios pagan sus facturas con un año de anticipación.
  \item un exceso de inversión en micromedidores ultrasónicos.
  \item la indexación tarifaria mensual automática.
\end{opts}

\q{La 'Depreciación Regulatoria' incluida por SUNASS en la estructura de costos de la tarifa quinquenal tiene como fin financiero principal:}
\begin{opts}
  \item financiar el subsidio cruzado de los usuarios industriales.
  \item constituir el Fondo de Inversión y Reposición de Activos de la EPS, garantizando que cuando una tubería o bomba culmine su vida útil regulatoria existan recursos monetarios tangibles para reemplazarla.
  \item distribuir utilidades extraordinarias al alcalde provincial.
  \item generar un escudo tributario para no pagar impuesto a la renta.
\end{opts}

\q{En los contratos de concesión de Asociaciones Público-Privadas (APP) para Plantas de Tratamiento de Aguas Residuales (como la PTAR Titicaca), la retribución al concesionario privado se estructura típicamente mediante dos componentes:}
\begin{opts}
  \item el cobro de una comisión por corte y reconexión en toda la región.
  \item la entrega de terrenos agrícolas en propiedad.
  \item la Remuneración por Inversiones (RPI), que amortiza el CAPEX de la planta, y la Remuneración por Mantenimiento y Operación (RPMO), que cubre el OPEX sujeto al cumplimiento de estándares de calidad del efluente.
  \item un arancel a la exportación y un crédito fiscal de IGV.
\end{opts}

\q{En el financiamiento de grandes plantas desalinizadoras o PTAR bajo la modalidad de 'Project Finance' (Financiamiento de Proyecto sin recurso):}
\begin{opts}
  \item los patrocinadores garantizan el 100\% de la deuda con su patrimonio corporativo personal.
  \item el Ministerio de Vivienda asume la totalidad del riesgo geológico.
  \item la tasa de interés del préstamo es subsidiada a tasa cero por la SBS.
  \item el repago de la deuda bancaria depende exclusivamente de los flujos de caja operativos generados por el propio proyecto y los activos de la sociedad de propósito especial (SPE) actúan como única garantía colateral.
\end{opts}

\q{Si una EPS municipal registra pérdidas contables acumuladas sistemáticas y no paga impuesto a la renta corporativo:}
\begin{opts}
  \item el beneficio del escudo fiscal de la deuda ($T \times K_d$) se pierde o se reduce drásticamente, haciendo que el financiamiento mediante endeudamiento sea financieramente menos ventajoso.
  \item los bancos le otorgan préstamos a tasa libre de riesgo.
  \item su WACC regulatorio se vuelve automáticamente negativo.
  \item la empresa debe disolverse obligatoriamente al término del año fiscal.
\end{opts}

\q{Los fondos de reserva para Mecanismos de Retribución por Servicios Ecosistémicos (MERESE) recaudados obligatoriamente en la tarifa de agua por SUNASS deben ser colocados por las EPS en:}
\begin{opts}
  \item el presupuesto de festividades del municipio provincial.
  \item cuentas bancarias intangibles y separadas en entidades del sistema financiero supervisadas por la SBS, utilizables exclusivamente para la ejecución de proyectos de conservación de cuencas aprobados.
  \item la cuenta corriente general para pago de combustibles de camionetas de la empresa.
  \item fondos de cobertura de alto riesgo en criptomonedas.
\end{opts}

\q{Una EPS que contrata un crédito financiero internacional en dólares estadounidenses para construir una planta potabilizadora mientras recauda sus ingresos tarifarios íntegramente en soles peruanos asume un severo riesgo de:}
\begin{opts}
  \item disminución del consumo de agua de los hogares.
  \item arbitraje de tasas de interés de Paridad Descubierta.
  \item descalce cambiario (Currency Mismatch), donde una depreciación del sol incrementa automáticamente la carga del servicio de la deuda en soles sin un aumento correlativo en los ingresos regulados.
  \item pérdida de la titularidad legal sobre las tuberías.
\end{opts}

\q{Al diseñar una Planta de Tratamiento modular donde la Etapa 1 se construye hoy para 500 litros/segundo y existe la 'Opción Real de Expansión' para duplicar a 1000 litros/segundo dentro de 8 años si la población crece:}
\begin{opts}
  \item SUNASS sanciona a la empresa por no construir la planta completa el primer día.
  \item la tasa de descuento social se eleva al 20\%.
  \item el costo total del proyecto se duplica innecesariamente en el año 0.
  \item se reduce el riesgo de sobreinversión en capacidad ociosa no remunerada y se preserva flexibilidad estratégica, maximizando el valor presente ajustado del proyecto.
\end{opts}

\q{Cuando el Organismo Técnico de la Administración de los Servicios de Saneamiento (OTASS) interviene una EPS y la incorpora al Régimen de Apoyo Transitorio (RAT), uno de los indicadores financieros determinantes es:}
\begin{opts}
  \item la causal de insolvencia financiera y patrimonio neto negativo, donde las pérdidas acumuladas han consumido la totalidad del capital social impidiendo la operación autónoma.
  \item el cumplimiento del 100\% de las metas del plan maestro.
  \item el exceso de liquidez corriente en moneda extranjera.
  \item la privatización de las fuentes de captación.
\end{opts}

\q{En el mecanismo de financiamiento de 'Obras por Impuestos' (OxI - Ley N° 29230) aplicado a obras de saneamiento rural:}
\begin{opts}
  \item las obras se ejecutan sin expediente técnico ni supervisión.
  \item una empresa privada financia y ejecuta la obra pública con cargo a deducir el monto invertido contra su pago a cuenta del Impuesto a la Renta de tercera categoría mediante CIPRL.
  \item el municipio exonera del pago de arbitrios a toda la población.
  \item la EPS emite bonos de titulización con garantía del Banco de la Nación.
\end{opts}

\q{La omisión del registro de pasivos ambientales y contingencias judiciales por vertimientos no autorizados en el Balance General de una EPS genera:}
\begin{opts}
  \item la anulación automática de los costos de financiamiento bancario.
  \item la duplicación del margen bruto de ventas.
  \item una sobrestimación del Patrimonio Neto contable de la empresa y una imagen financiera artificialmente saneada que oculta riesgos latentes de multas de la OEFA y SUNASS.
  \item un incremento inmediato de la recaudación tarifaria.
\end{opts}

\q{Al comparar el reemplazo de micromedidores mecánicos convencionales (costo inicial bajo, vida útil de 5 años, submedición por desgaste rápido) frente a medidores ultrasónicos estáticos (costo inicial mayor, vida útil de 12 años, precisión sostenida y telemetría):}
\begin{opts}
  \item los medidores ultrasónicos no pueden ser depreciados tributariamente.
  \item la teoría de las finanzas exige no medir el consumo de agua.
  \item se debe elegir siempre el medidor mecánico por tener menor costo en $t=0$.
  \item el criterio de evaluación financiera idóneo es el Costo Anual Equivalente (CAE), incorporando tanto los costos de reposición como los ingresos adicionales recuperados por reducción de subfacturación.
\end{opts}

\q{El periodo regulatorio tarifario fijado por SUNASS tiene una vigencia estándar de cinco (5) años (quinquenio regulatorio). Esta duración temporal se fundamenta financieramente en:}
\begin{opts}
  \item el ciclo de maduración y ejecución de las inversiones de capital (CAPEX) y la necesidad de otorgar un marco de estabilidad y certidumbre para la recuperación del WACC antes de transferir las ganancias de productividad a los usuarios.
  \item la rotación anual del directorio de la EPS.
  \item el periodo de mandato de los regidores municipales.
  \item la frecuencia del censo nacional del INEI.
\end{opts}

\q{En un contrato de concesión de Asociación Público-Privada (APP) de saneamiento a 25 años, una cláusula de 'Restablecimiento del Equilibrio Económico-Financiero' protege al inversionista privado y al Estado ante:}
\begin{opts}
  \item cambios ordinarios en los precios de mercado que forman parte del riesgo comercial propio del negocio.
  \item eventos sobrevinientes y extraordinarios no atribuibles a las partes, tales como modificaciones normativas regulatorias imprevistas (cambio en la ley), catástrofes naturales de fuerza mayor o variaciones fiscales drásticas.
  \item la decisión unilateral del concesionario de suspender el tratamiento de aguas.
  \item la reducción voluntaria de la jornada laboral de los obreros de la planta.
\end{opts}


\newpage
\section*{Tabla de claves rápidas}

\noindent Use esta tabla para una verificación inmediata (preguntas 1 a 180). El solucionario con la justificación de cada respuesta se presenta en la sección siguiente.

\vspace{0.35cm}
\begin{center}
\footnotesize
\renewcommand{\arraystretch}{1.25}
\begin{tabular}{|c|c|c|c|c|c|c|c|c|c|c|}
\hline
\textbf{N.} & \textbf{1} & \textbf{2} & \textbf{3} & \textbf{4} & \textbf{5} & \textbf{6} & \textbf{7} & \textbf{8} & \textbf{9} & \textbf{10} \\
\hline
\textbf{Clave} & C & A & C & B & D & C & B & C & A & A \\
\hline
\textbf{N.} & \textbf{11} & \textbf{12} & \textbf{13} & \textbf{14} & \textbf{15} & \textbf{16} & \textbf{17} & \textbf{18} & \textbf{19} & \textbf{20} \\
\hline
\textbf{Clave} & B & A & B & D & B & A & C & A & D & B \\
\hline
\textbf{N.} & \textbf{21} & \textbf{22} & \textbf{23} & \textbf{24} & \textbf{25} & \textbf{26} & \textbf{27} & \textbf{28} & \textbf{29} & \textbf{30} \\
\hline
\textbf{Clave} & A & B & D & D & B & B & A & B & C & C \\
\hline
\textbf{N.} & \textbf{31} & \textbf{32} & \textbf{33} & \textbf{34} & \textbf{35} & \textbf{36} & \textbf{37} & \textbf{38} & \textbf{39} & \textbf{40} \\
\hline
\textbf{Clave} & D & D & A & B & B & C & A & A & C & B \\
\hline
\textbf{N.} & \textbf{41} & \textbf{42} & \textbf{43} & \textbf{44} & \textbf{45} & \textbf{46} & \textbf{47} & \textbf{48} & \textbf{49} & \textbf{50} \\
\hline
\textbf{Clave} & C & B & B & C & C & D & B & A & C & B \\
\hline
\textbf{N.} & \textbf{51} & \textbf{52} & \textbf{53} & \textbf{54} & \textbf{55} & \textbf{56} & \textbf{57} & \textbf{58} & \textbf{59} & \textbf{60} \\
\hline
\textbf{Clave} & A & D & B & C & B & B & D & D & C & D \\
\hline
\textbf{N.} & \textbf{61} & \textbf{62} & \textbf{63} & \textbf{64} & \textbf{65} & \textbf{66} & \textbf{67} & \textbf{68} & \textbf{69} & \textbf{70} \\
\hline
\textbf{Clave} & D & C & D & A & A & C & D & B & C & A \\
\hline
\textbf{N.} & \textbf{71} & \textbf{72} & \textbf{73} & \textbf{74} & \textbf{75} & \textbf{76} & \textbf{77} & \textbf{78} & \textbf{79} & \textbf{80} \\
\hline
\textbf{Clave} & C & A & C & D & B & D & D & C & D & A \\
\hline
\textbf{N.} & \textbf{81} & \textbf{82} & \textbf{83} & \textbf{84} & \textbf{85} & \textbf{86} & \textbf{87} & \textbf{88} & \textbf{89} & \textbf{90} \\
\hline
\textbf{Clave} & A & D & A & A & D & B & A & A & C & D \\
\hline
\textbf{N.} & \textbf{91} & \textbf{92} & \textbf{93} & \textbf{94} & \textbf{95} & \textbf{96} & \textbf{97} & \textbf{98} & \textbf{99} & \textbf{100} \\
\hline
\textbf{Clave} & B & C & D & A & C & D & A & B & D & A \\
\hline
\textbf{N.} & \textbf{101} & \textbf{102} & \textbf{103} & \textbf{104} & \textbf{105} & \textbf{106} & \textbf{107} & \textbf{108} & \textbf{109} & \textbf{110} \\
\hline
\textbf{Clave} & B & C & A & B & C & D & B & C & D & A \\
\hline
\textbf{N.} & \textbf{111} & \textbf{112} & \textbf{113} & \textbf{114} & \textbf{115} & \textbf{116} & \textbf{117} & \textbf{118} & \textbf{119} & \textbf{120} \\
\hline
\textbf{Clave} & C & D & A & B & D & A & B & C & A & B \\
\hline
\textbf{N.} & \textbf{121} & \textbf{122} & \textbf{123} & \textbf{124} & \textbf{125} & \textbf{126} & \textbf{127} & \textbf{128} & \textbf{129} & \textbf{130} \\
\hline
\textbf{Clave} & C & D & B & C & D & A & C & D & A & B \\
\hline
\textbf{N.} & \textbf{131} & \textbf{132} & \textbf{133} & \textbf{134} & \textbf{135} & \textbf{136} & \textbf{137} & \textbf{138} & \textbf{139} & \textbf{140} \\
\hline
\textbf{Clave} & D & A & B & C & A & B & C & D & B & C \\
\hline
\textbf{N.} & \textbf{141} & \textbf{142} & \textbf{143} & \textbf{144} & \textbf{145} & \textbf{146} & \textbf{147} & \textbf{148} & \textbf{149} & \textbf{150} \\
\hline
\textbf{Clave} & D & A & C & D & A & B & D & A & B & C \\
\hline
\textbf{N.} & \textbf{151} & \textbf{152} & \textbf{153} & \textbf{154} & \textbf{155} & \textbf{156} & \textbf{157} & \textbf{158} & \textbf{159} & \textbf{160} \\
\hline
\textbf{Clave} & A & B & C & D & B & C & D & A & C & D \\
\hline
\textbf{N.} & \textbf{161} & \textbf{162} & \textbf{163} & \textbf{164} & \textbf{165} & \textbf{166} & \textbf{167} & \textbf{168} & \textbf{169} & \textbf{170} \\
\hline
\textbf{Clave} & A & B & D & A & B & C & A & B & C & D \\
\hline
\textbf{N.} & \textbf{171} & \textbf{172} & \textbf{173} & \textbf{174} & \textbf{175} & \textbf{176} & \textbf{177} & \textbf{178} & \textbf{179} & \textbf{180} \\
\hline
\textbf{Clave} & B & C & D & A & C & D & A & B & D & A \\
\hline
\end{tabular}
\end{center}

\newpage
\section*{Solucionario completo}

\noindent Cada respuesta incluye la clave oficial y su justificación analítica o matemática detallada.

\subsection*{Bloque A. Fundamentos teóricos y conceptuales}

\sol{1}{C}{\textbf{Tema:} Valor del dinero en el tiempo. \textbf{Explicación:} El dinero tiene valor temporal por tres factores intrínsecos: la impaciencia racional o preferencia por consumir hoy, la pérdida de poder de compra por inflación y el costo de oportunidad del capital (posibilidad de generar intereses libres de riesgo).}

\sol{2}{A}{\textbf{Tema:} Interés compuesto vs simple. \textbf{Explicación:} En el interés simple el rendimiento se calcula siempre sobre el principal inicial ($VF = VP(1 + i \cdot n)$). En el interés compuesto los intereses devengados se reinvierten sistemáticamente, generando un crecimiento exponencial ($VF = VP(1 + i)^n$).}

\sol{3}{C}{\textbf{Tema:} Tasa Efectiva Anual (TEA). \textbf{Explicación:} La TNA es una tasa de referencia meramente proporcional. La TEA refleja el rendimiento o costo efectivo generado al cabo de un año cuando los intereses se capitalizan $m$ veces al año (mensual, diario, etc.), cumpliéndose siempre $TEA > TNA$ para capitalización subanual.}

\sol{4}{B}{\textbf{Tema:} Ecuación de Fisher exacta. \textbf{Explicación:} La versión exacta de Irving Fisher descuenta el efecto inflacionario del principal y de los intereses: $(1 + i) = (1 + r)(1 + \pi) = 1 + r + \pi + r\pi$. La aproximación lineal $i \approx r + \pi$ ignora el término cruzado $r\pi$, el cual es relevante cuando la inflación es moderada o alta.}

\sol{5}{D}{\textbf{Tema:} Valor presente de una perpetuidad simple. \textbf{Explicación:} Sumando los términos de una serie geométrica infinita decreciente con razón $\frac{1}{1+r}$: $VP = \sum_{t=1}^{\infty} \frac{C}{(1+r)^t} = \frac{C}{1+r} \frac{1}{1 - 1/(1+r)} = \frac{C}{r}$.}

\sol{6}{C}{\textbf{Tema:} Modelo de Gordon-Shapiro. \textbf{Explicación:} Cuando los dividendos o flujos crecen geométricamente a tasa $g$, el valor presente de la serie infinita converge formalmente a $VP = \frac{C_1}{r - g}$, siempre y cuando la tasa de descuento supere a la tasa de crecimiento ($r > g$).}

\sol{7}{B}{\textbf{Tema:} Anualidades anticipadas vs vencidas. \textbf{Explicación:} Al recibirse las cuotas al inicio de cada periodo, cada flujo está descontado por un factor $(1+r)$ menos que en la anualidad ordinaria. Por ende, $VP_{ant} = VP_{venc} \times (1 + r)$.}

\sol{8}{C}{\textbf{Tema:} Pendiente de la curva de rendimiento. \textbf{Explicación:} La teoría de la prima de liquidez (Hicks) explica que ante el mayor riesgo de duración y volatilidad del precio de los activos a largo plazo, los agentes demandan un mayor retorno esperado para inmovilizar su liquidez en horizontes dilatados.}

\sol{9}{A}{\textbf{Tema:} Ecuación fundamental de la contabilidad. \textbf{Explicación:} El Activo Total (totalidad de bienes y derechos controlados por la entidad) está íntegramente financiado por dos fuentes: deuda con terceros ajenos a la empresa (Pasivo) y capital aportado o retenido por los propietarios (Patrimonio Neto).}

\sol{10}{A}{\textbf{Tema:} Prueba ácida vs Razón corriente. \textbf{Explicación:} $Prueba\ \acute{A}cida = \frac{Activo\ Corriente - Inventarios}{Pasivo\ Corriente}$. Permite evaluar si la empresa puede honrar sus deudas exigibles de corto plazo sin depender de la venta o liquidación forzosa de sus existencias.}

\sol{11}{B}{\textbf{Tema:} Capital de Trabajo Neto. \textbf{Explicación:} El CTN mide el margen de seguridad financiera de corto plazo: la porción de los activos circulantes financiada con pasivos de largo plazo o recursos patrimoniales permanentes.}

\sol{12}{A}{\textbf{Tema:} Ratio de apalancamiento D/E. \textbf{Explicación:} El cociente $D/E = Pasivo\ Total / Patrimonio\ Neto$ cuantifica el riesgo financiero: un elevado apalancamiento amplifica los retornos del accionista en periodos de bonanza pero eleva drásticamente la probabilidad de quiebra ante caídas en los ingresos operativos.}

\sol{13}{B}{\textbf{Tema:} Ratio de cobertura de intereses. \textbf{Explicación:} Indica cuántas veces el beneficio generado por las operaciones del negocio cubre la carga fija de intereses de la deuda. Ratios inferiores a 1.5 a 2.0 veces alertan a los bancos sobre alto riesgo de default crediticio.}

\sol{14}{D}{\textbf{Tema:} ROA vs ROE. \textbf{Explicación:} $ROA = \frac{Utilidad\ Neta}{Activo\ Total}$ mide la productividad integral de las inversiones; $ROE = \frac{Utilidad\ Neta}{Patrimonio}$ mide el rendimiento financiero del propietario, el cual se apalanca con el uso de endeudamiento si el costo de la deuda es menor que el retorno de los activos.}

\sol{15}{B}{\textbf{Tema:} Análisis DuPont del ROE. \textbf{Explicación:} $ROE = \frac{Utilidad\ Neta}{Ventas} \times \frac{Ventas}{Activos} \times \frac{Activos}{Patrimonio} = Margen \times Rotaci\acute{o}n \times Apalancamiento$. Permite identificar si la rentabilidad proviene de margen de precios, volumen de ventas o endeudamiento.}

\sol{16}{A}{\textbf{Tema:} Ciclo de conversión de efectivo. \textbf{Explicación:} $CCE = D\acute{\imath}as\ Inventario + D\acute{\imath}as\ Cobranza - D\acute{\imath}as\ Pago$. Un menor CCE reduce los requerimientos de financiamiento de capital de trabajo y eleva la liquidez del negocio.}

\sol{17}{C}{\textbf{Tema:} Flujo de caja libre del proyecto. \textbf{Explicación:} El FCLP evalúa la rentabilidad económica del activo puro. Los intereses y amortizaciones no se restan del flujo porque el costo de financiamiento de la deuda ya está incorporado en la tasa de descuento utilizada para descontarlo ($WACC$).}

\sol{18}{A}{\textbf{Tema:} Tratamiento de costos hundidos. \textbf{Explicación:} La regla de decisión financiera exige evaluar únicamente flujos de caja incrementales futuros ($Marginales$). Los costos pasados e irreversibles son irrelevantes para la decisión prospectiva de seguir o abandonar el proyecto.}

\sol{19}{D}{\textbf{Tema:} Regla de decisión del VAN. \textbf{Explicación:} El $VAN = \sum_{t=1}^n \frac{FC_t}{(1+r)^t} - I_0$ cuantifica el incremento absoluto en la riqueza de la empresa por encima del costo de oportunidad del capital. Todo proyecto con $VAN > 0$ genera valor y debe ejecutarse.}

\sol{20}{B}{\textbf{Tema:} Aditividad de valor del VAN. \textbf{Explicación:} La aditividad de valor permite evaluar proyectos de manera descentralizada o sumar el valor de divisiones de negocio sin distorsiones matemáticas, propiedad que no cumple la TIR ni el Payback.}

\sol{21}{A}{\textbf{Tema:} Definición matemática de la TIR. \textbf{Explicación:} La TIR es la solución $r$ a la ecuación polinómica $\sum_{t=0}^n \frac{FC_t}{(1+r)^t} = 0$. Representa la rentabilidad porcentual interna intrínseca que rinde el capital invertido en el proyecto a lo largo de su vida útil.}

\sol{22}{B}{\textbf{Tema:} Supuesto de reinversión VAN vs TIR. \textbf{Explicación:} Si un proyecto tiene una TIR del 80\%, la fórmula asume que cada flujo recibido puede reinvertirse de inmediato al 80\%, inflando artificialmente el atractivo del proyecto. El VAN es más realista al asumir reinversión al costo de oportunidad de mercado ($WACC$).}

\sol{23}{D}{\textbf{Tema:} TIR múltiple por cambios de signo. \textbf{Explicación:} Cada cambio de signo en el polinomio de flujos de caja genera una raíz real positiva potencial. Al haber múltiples TIR, no existe una tasa única con la cual comparar el costo de capital, debiendo descartarse la TIR a favor del VAN.}

\sol{24}{D}{\textbf{Tema:} Conflicto VAN vs TIR en proyectos excluyentes. \textbf{Explicación:} El objetivo de la corporación es maximizar la riqueza total de los accionistas (valor absoluto en soles), no la tasa porcentual por sol invertido. Ante conflicto entre proyectos excluyentes, la regla del VAN superior manda siempre sobre la TIR.}

\sol{25}{B}{\textbf{Tema:} Ventajas de la TIRM. \textbf{Explicación:} La TIRM combina todos los flujos de entrada en un único Valor Futuro acumulado a la tasa de costo de capital y los desembolsos en un único Valor Presente. Al haber un solo cambio de signo temporal, garantiza una tasa única y sin el sesgo de reinversión a tasas extremas.}

\sol{26}{B}{\textbf{Tema:} Limitaciones del Payback simple. \textbf{Explicación:} El Payback suma flujos nominales sin descontar (tratando un sol del año 5 igual que un sol de hoy) y castiga arbitrariamente proyectos estratégicos de larga vida útil que generan grandes beneficios tras el año de corte establecido.}

\sol{27}{A}{\textbf{Tema:} Costo Anual Equivalente (CAE). \textbf{Explicación:} Cuando dos plantas o bombas de agua tienen distinta duración (ej. 5 años vs 10 años) y deben reponerse indefinidamente, no se pueden comparar sus VAN directos. El CAE anualiza el costo presente total mediante una anualidad uniforme ($CAE = VP(Costos) / A_n^r$), eligiéndose el menor CAE.}

\sol{28}{B}{\textbf{Tema:} Fórmula del WACC corporativo. \textbf{Explicación:} El WACC pondera el costo de cada fuente de financiamiento por su participación en el valor total de la firma ($V$). La deuda incorpora el factor $(1 - T)$ porque el pago de intereses es deducible del impuesto a la renta (escudo fiscal).}

\sol{29}{C}{\textbf{Tema:} Concepto de Beta en el CAPM. \textbf{Explicación:} El inversionista en portafolios diversificados elimina el riesgo no sistemático (idiosincrático). Por tanto, el mercado solo premia el riesgo que no se puede diversificar (riesgo sistemático medido por $\beta = \frac{Cov(R_i, R_m)}{Var(R_m)}$).}

\sol{30}{C}{\textbf{Tema:} Modigliani-Miller Proposición I (sin impuestos). \textbf{Explicación:} Bajo ausencia de impuestos y fricciones, el apalancamiento no genera valor: el mayor riesgo financiero de la deuda eleva el costo de los recursos propios ($K_e$) en la proporción exacta para mantener constante el $WACC$, dejando el valor de la firma inalterado.}

\subsection*{Bloque B. Aplicaciones cuantitativas y financieras}

\sol{31}{D}{\textbf{Tema:} Conversión de TNA mensual a TEA. \textbf{Explicación:} La tasa efectiva mensual es $i_m = 12\% / 12 = 1\% = 0.01$. Capitalizando a 12 meses: $TEA = (1 + 0.01)^{12} - 1 = (1.01)^{12} - 1 \approx 1.126825 - 1 = 12.68\%$.}

\sol{32}{D}{\textbf{Tema:} Cálculo de tasa real de Fisher. \textbf{Explicación:} Aplicando la fórmula exacta: $1 + r = \frac{1 + i}{1 + \pi} = \frac{1.0815}{1.0300} = 1.0500 \implies r = 5.00\%$ anual.}

\sol{33}{A}{\textbf{Tema:} Valor presente de perpetuidad simple. \textbf{Explicación:} $VP = \frac{C}{r} = \frac{50000}{0.08} = 625000$ soles.}

\sol{34}{B}{\textbf{Tema:} Cálculo con modelo de Gordon-Shapiro. \textbf{Explicación:} $VP = \frac{C_1}{r - g} = \frac{30000}{0.09 - 0.03} = \frac{30000}{0.06} = 500000$ soles.}

\sol{35}{B}{\textbf{Tema:} Cálculo de cuota fija de anualidad. \textbf{Explicación:} $C = P \times \left[ \frac{r(1+r)^n}{(1+r)^n - 1} \right] = 100000 \times 0.2637975 = 26379.75$ soles anuales.}

\sol{36}{C}{\textbf{Tema:} Cálculo de razón corriente y prueba ácida. \textbf{Explicación:} $Razón\ Corriente = \frac{160000}{80000} = 2.0$ veces. $Prueba\ \acute{A}cida = \frac{160000 - 60000}{80000} = \frac{100000}{80000} = 1.25$ veces.}

\sol{37}{A}{\textbf{Tema:} Cálculo de D/E ratio. \textbf{Explicación:} Patrimonio $= Activo - Pasivo = 500000 - 300000 = 200000$ soles. $D/E = \frac{Deuda}{Patrimonio} = \frac{300000}{200000} = 1.50$ (o $150\%$).}

\sol{38}{A}{\textbf{Tema:} Cálculo de cobertura de intereses. \textbf{Explicación:} $Cobertura = \frac{EBIT}{Intereses} = \frac{120000}{30000} = 4.00$ veces. La empresa genera 4 soles de beneficio operativo por cada sol que debe pagar en intereses.}

\sol{39}{C}{\textbf{Tema:} Cálculo del ROE con DuPont. \textbf{Explicación:} $ROE = Margen \times Rotaci\acute{o}n \times Apalancamiento = 0.08 \times 1.5 \times 1.50 = 0.18 = 18.0\%$.}

\sol{40}{B}{\textbf{Tema:} Cálculo del ciclo de conversión del efectivo. \textbf{Explicación:} $CCE = D\acute{\imath}as\ Cobro + D\acute{\imath}as\ Inventario - D\acute{\imath}as\ Pago = 45 + 60 - 40 = 65$ días.}

\sol{41}{C}{\textbf{Tema:} Cálculo de Flujo de Caja Libre (FCLP). \textbf{Explicación:} $FCLP = EBIT(1 - T) + Depreciaci\acute{o}n - CAPEX - \Delta CTN = 200000(0.70) + 40000 - 50000 - 10000 = 140000 + 40000 - 60000 = 120000$ soles.}

\sol{42}{B}{\textbf{Tema:} Cálculo del VAN de un proyecto. \textbf{Explicación:} $VAN = -100000 + \frac{40000}{1.10} + \frac{40000}{(1.10)^2} + \frac{40000}{(1.10)^3} = -100000 + 36363.6 + 33057.9 + 30052.6 = -525.9$ soles. Al ser $VAN < 0$, el proyecto debe rechazarse.}

\sol{43}{B}{\textbf{Tema:} Cálculo de TIR de 2 periodos. \textbf{Explicación:} Sea $x = 1/(1+TIR)$. La ecuación es $10000 = 6000x + 6000x^2 \implies 6x^2 + 6x - 10 = 0 \implies 3x^2 + 3x - 5 = 0$. $x = \frac{-3 + \sqrt{9 - 4(3)(-5)}}{6} = \frac{-3 + \sqrt{69}}{6} \approx \frac{-3 + 8.3066}{6} = 0.88443$. Luego $1+TIR = 1/0.88443 \approx 1.13066 \implies TIR \approx 13.07\%$.}

\sol{44}{C}{\textbf{Tema:} Cálculo de la Tasa de Fisher. \textbf{Explicación:} La diferencia de flujos (A - B) es: $t_0 = 0$, $t_1 = 700 - 200 = 500$, $t_2 = 700 - 1300 = -600$. Igualando a cero: $\frac{500}{1+r} - \frac{600}{(1+r)^2} = 0 \implies 500 = \frac{600}{1+r} \implies 1 + r = \frac{600}{500} = 1.20 \implies r_{Fisher} = 20.0\%$.}

\sol{45}{C}{\textbf{Tema:} Cálculo de Payback simple. \textbf{Explicación:} $Payback = \frac{Inversi\acute{o}n}{Flujo\ Anual} = \frac{70000}{25000} = 2.80$ años.}

\sol{46}{D}{\textbf{Tema:} Relación Beneficio/Costo (B/C). \textbf{Explicación:} $B/C = \frac{VP(Beneficios)}{VP(Costos)} = \frac{280000}{200000} = 1.40$. Al ser $B/C > 1$, el proyecto es rentable y genera un $VAN = 280000 - 200000 = +80000$ soles.}

\sol{47}{B}{\textbf{Tema:} Cálculo de Costo Anual Equivalente (CAE). \textbf{Explicación:} $CAE_{Alfa} = \frac{50000}{2.48685} \approx 20105.75$ soles. $CAE_{Beta} = \frac{75000}{3.79079} \approx 19784.80$ soles. Aunque Beta requiere mayor costo inicial total, su mayor durabilidad genera un menor costo anualizado ($19784.8 < 20105.7$).}

\sol{48}{A}{\textbf{Tema:} Cálculo de WACC con impuestos. \textbf{Explicación:} $WACC = K_e (E/V) + K_d (1 - T)(D/V) = 0.12(0.60) + 0.07(1 - 0.30)(0.40) = 0.072 + 0.07(0.70)(0.40) = 0.072 + 0.0196 = 0.0916 = 9.16\%$.}

\sol{49}{C}{\textbf{Tema:} Cálculo de costo de capital con CAPM emergente. \textbf{Explicación:} $K_e = R_f + \beta(R_m - R_f) + RP = 0.04 + 0.75(0.06) + 0.015 = 0.04 + 0.045 + 0.015 = 0.1000 = 10.00\%$.}

\sol{50}{B}{\textbf{Tema:} Cálculo del escudo fiscal de la deuda. \textbf{Explicación:} El ahorro tributario anual en intereses es $T \times (K_d \times D) = 0.30 \times 0.06 \times 3000000 = 54000$ soles. Descontando esta perpetuidad al costo de la deuda $K_d$: $VP = \frac{T \times K_d \times D}{K_d} = T \times D = 0.30 \times 3000000 = 900000$ soles.}

\sol{51}{A}{\textbf{Tema:} Desglose de cuota en sistema francés. \textbf{Explicación:} El interés del primer periodo se calcula sobre el saldo inicial: $I_1 = 100000 \times 0.10 = 10000.00$ soles. La amortización es la diferencia: $A_1 = Cuota - I_1 = 31547.08 - 10000.00 = 21547.08$ soles.}

\sol{52}{D}{\textbf{Tema:} Sistema de amortización alemán. \textbf{Explicación:} Amortización fija $= 60000/3 = 20000$ soles. Año 1: interés $6000$, cuota $26000$. Año 2: saldo $40000$, interés $4000$, cuota $24000$. Año 3: saldo $20000$, interés $2000$, cuota $22000$.}

\sol{53}{B}{\textbf{Tema:} Precio de bono cupón cero. \textbf{Explicación:} $P = \frac{Valor\ Nominal}{(1 + YTM)^n} = \frac{1000}{(1.08)^2} \approx 857.34$ soles.}

\sol{54}{C}{\textbf{Tema:} Sensibilidad de precios por Duración Modificada. \textbf{Explicación:} $\frac{\Delta P}{P} \approx -D_{mod} \times \Delta y = -4.5 \times (+0.005) = -0.0225 = -2.25\%$. El precio del bono cae un $2.25\%$.}

\sol{55}{B}{\textbf{Tema:} Cálculo de la TIRM. \textbf{Explicación:} $1000(1 + TIRM)^2 = 1460 \implies (1 + TIRM)^2 = 1.46 \implies 1 + TIRM = 1.2083 \implies TIRM = 20.83\%$.}

\sol{56}{B}{\textbf{Tema:} Punto de equilibrio operativo. \textbf{Explicación:} $Q = \frac{Costos\ Fijos}{P - CV} = \frac{120000}{5.0 - 2.0} = \frac{120000}{3.0} = 40000$ metros cúbicos.}

\sol{57}{D}{\textbf{Tema:} Grado de apalancamiento operativo (GAO). \textbf{Explicación:} $\% \Delta EBIT = GAO \times \% \Delta Ventas = 2.5 \times 8.0\% = 20.0\%$.}

\sol{58}{D}{\textbf{Tema:} Grado de apalancamiento total (GAT). \textbf{Explicación:} $GAT = GAO \times GAF = 2.0 \times 1.5 = 3.0$. Mide el impacto porcentual total sobre la utilidad neta por acción ante una variación del 1\% en las ventas.}

\sol{59}{C}{\textbf{Tema:} Método del Valor Presente Ajustado (APV). \textbf{Explicación:} $APV = VAN_{unlevered} + VP(Financiamiento)$. Desagrega el valor intrínseco del proyecto de los efectos colaterales del financiamiento, siendo ideal cuando el nivel de deuda no es una proporción constante del valor de mercado.}

\sol{60}{D}{\textbf{Tema:} Opciones reales de expansión en infraestructura. \textbf{Explicación:} La flexibilidad gerencial para ampliar capacidad si las condiciones de demanda son favorables tiene valor económico real (Opción Call). El enfoque de opciones reales reconoce que la gerencia no está atada pasivamente a un plan rígido inicial.}

\sol{61}{D}{\textbf{Tema:} Modigliani-Miller Proposición II con impuestos. \textbf{Explicación:} $K_e = K_u + (K_u - K_d)(1 - T)(D/E) = 0.10 + (0.10 - 0.06)(0.70)(0.50) = 0.10 + 0.04(0.35) = 0.10 + 0.014 = 0.114 = 11.40\%$.}

\sol{62}{C}{\textbf{Tema:} Consistencia en evaluación con inflación. \textbf{Explicación:} Mezclar flujos nominales con tasas reales (o viceversa) introduce severas distorsiones. La consistencia matemática exige que el numerador y el denominador incorporen (o excluyan) la prima de inflación en idéntica medida.}

\sol{63}{D}{\textbf{Tema:} Valor terminal por perpetuidad. \textbf{Explicación:} $Valor\ Terminal_{10} = \frac{FCLP_{10}}{r} = \frac{80000}{0.10} = 800000$ soles evaluados en el año 10. Para sumarlo al VAN se debe descontar a valor presente: $\frac{800000}{(1.10)^{10}}$.}

\sol{64}{A}{\textbf{Tema:} Cálculo de ratio P/E. \textbf{Explicación:} $P/E = \frac{Precio}{Utilidad\ por\ Acci\acute{o}n} = \frac{30.0}{2.50} = 12.0$ veces. Los inversionistas pagan 12 soles por cada sol de beneficio anual generado por la compañía.}

\sol{65}{A}{\textbf{Tema:} Cálculo de YTM de bono cupón cero. \textbf{Explicación:} $(1 + YTM)^3 = \frac{1000}{751.3146} = 1.331 \implies 1 + YTM = \sqrt[3]{1.331} = 1.10 \implies YTM = 10.00\%$ anual.}

\subsection*{Bloque C. Casos con implicancias regulatorias y sector saneamiento}

\sol{66}{C}{\textbf{Tema:} Estructura de capital regulatoria de SUNASS. \textbf{Explicación:} SUNASS establece una estructura de financiamiento eficiente y normativa (empresa modelo eficiente) para que la tarifa remunere un costo de capital prudente, sin convalidar sobrecostos si la EPS está mal financiada ni desincentivar la inversión si carece de deuda.}

\sol{67}{D}{\textbf{Tema:} Tasa libre de riesgo en regulación. \textbf{Explicación:} Los US Treasuries a largo plazo son el estándar financiero internacional para $R_f$. Dado que la infraestructura de agua tiene vidas útiles de 20 a 40 años, se eligen bonos soberanos con vencimientos largos compatibles con la vida de los activos.}

\sol{68}{B}{\textbf{Tema:} Riesgo país EMBI+ en WACC regulatorio. \textbf{Explicación:} El EMBI+ cuantifica el riesgo soberano de impago y estabilidad macroeconómica de invertir en el Perú frente a EE.UU., sumándose a la tasa libre de riesgo en el modelo CAPM extendido para mercados emergentes.}

\sol{69}{C}{\textbf{Tema:} Definición de Base de Activos Regulatorios. \textbf{Explicación:} La BAR representa el capital físico productivo reconocido por el regulador (plantas, redes, reservorios). La retribución al capital en la tarifa se calcula como $Retribuci\acute{o}n = BAR \times WACC$, asegurando el retorno del capital invertido.}

\sol{70}{A}{\textbf{Tema:} Metodología del Valor Nuevo de Reemplazo (VNR). \textbf{Explicación:} El VNR simula el costo de reproducir el servicio hoy con la mejor tecnología disponible y eficiente. Evita subvaluar activos amortizados contablemente que aún operan y evita sobrevaluar infraestructura obsoleta con tecnología costosa superada.}

\sol{71}{C}{\textbf{Tema:} Justificación del WACC en empresas públicas. \textbf{Explicación:} Si la tarifa solo cubriera costos operativos (OPEX) y depreciación contable con retorno cero, la empresa descapitalizaría sus redes y sería incapaz de financiar nuevas plantas de tratamiento o acceder a préstamos del BID o KfW sin depender de subsidios del MEF.}

\sol{72}{A}{\textbf{Tema:} Evaluación social en Invierte.pe. \textbf{Explicación:} Invierte.pe evalúa si el proyecto genera bienestar social neto para el país. Corrige distorsiones de mercado (impuestos indirectos, subsidios, desempleo) mediante factores de conversión social aplicados a la mano de obra, divisa y combustibles.}

\sol{73}{C}{\textbf{Tema:} Tasa Social de Descuento del MEF. \textbf{Explicación:} La TSD oficial de Invierte.pe se ubica en $8.0\%$ anual real (reducida desde el histórico 10\% y 9\%). Representa el costo de oportunidad social de los fondos públicos utilizados en proyectos de inversión pública en el Perú.}

\sol{74}{D}{\textbf{Tema:} Factor de corrección de mano de obra no calificada. \textbf{Explicación:} En regiones con subempleo, emplear a un trabajador rural en una obra de saneamiento no reduce la producción agrícola en la misma cuantía del salario pagado. El factor de corrección descuenta esa brecha para reflejar su menor costo social real.}

\sol{75}{B}{\textbf{Tema:} Análisis Costo-Eficacia en Invierte.pe. \textbf{Explicación:} El Análisis Costo-Eficacia se utiliza cuando los beneficios sociales son evidentes pero no monetizables con rigor econométrico. Se fija el objetivo social (ej. 100\% de cobertura y agua clorada) y se elige la alternativa con menor Ratio Costo-Eficacia ($RCE$).}

\sol{76}{D}{\textbf{Tema:} Ratio DSCR en contratos de crédito de EPS. \textbf{Explicación:} El DSCR (Debt Service Coverage Ratio) mide la holgura de caja de la empresa prestadora para pagar amortización e intereses con el flujo generado por las tarifas. Si cae por debajo de 1.20, los bancos restringen nuevos desembolsos.}

\sol{77}{D}{\textbf{Tema:} Cuentas por cobrar del sector público en EPS. \textbf{Explicación:} La morosidad de entidades del Estado con las EPS genera un grave problema de capital de trabajo: las EPS deben pagar energía eléctrica, cloro y planillas al contado, mientras sus cuentas por cobrar públicas tardan meses en liquidarse.}

\sol{78}{C}{\textbf{Tema:} Depreciación regulatoria y fondo de reposición. \textbf{Explicación:} A diferencia de una empresa comercial donde la depreciación es un mero asiento contable de costo hundido, en la regulación de SUNASS la depreciación alimenta cuentas bancarias restringidas de inversión para garantizar la reposición continua de la infraestructura.}

\sol{79}{D}{\textbf{Tema:} Estructura RPI y RPMO en concesiones de saneamiento. \textbf{Explicación:} El esquema RPI/RPMO separa el repago del capital invertido (compromiso firme de pago del Estado o de la tarifa) del costo operativo mensual, el cual se penaliza si la planta no cumple con los límites máximos permisibles de vertimiento.}

\sol{80}{A}{\textbf{Tema:} Project Finance en infraestructura hídrica. \textbf{Explicación:} En Project Finance los prestamistas no tienen recurso contra los accionistas promotores si el proyecto quiebra: el riesgo descansa en la viabilidad intrínseca del flujo de caja del contrato de concesión y en la solidez del marco regulatorio tarifario.}

\sol{81}{A}{\textbf{Tema:} Pérdida del escudo fiscal en empresas insolventes. \textbf{Explicación:} El escudo fiscal de los intereses solo genera valor en efectivo si la empresa tiene utilidades tributarias gravables contra las cuales deducir el gasto. Sin ganancias operativas netas, el beneficio del escudo impositivo de Modigliani-Miller desaparece.}

\sol{82}{D}{\textbf{Tema:} Intangibilidad de fondos MERESE. \textbf{Explicación:} La regulación de SUNASS exige estricta intangibilidad contable y financiera: los fondos de la cuota MERESE y GRD no pueden mezclarse con el flujo operativo corriente de la EPS ni destinarse a gastos administrativos generales.}

\sol{83}{A}{\textbf{Tema:} Descalce cambiario en empresas de servicios públicos. \textbf{Explicación:} Como las tarifas de agua no están dolarizadas, un salto en el tipo de cambio encarece el pago de amortización e intereses en moneda nacional, comprometiendo la liquidez y solvencia de la prestadora si no cuenta con cobertura cambiaria (hedging).}

\sol{84}{A}{\textbf{Tema:} Opciones reales modulares en saneamiento. \textbf{Explicación:} La construcción modular es una opción real de crecimiento: evita hundir capital en megaobras con capacidad ociosa durante décadas y permite adaptar la inversión según la evolución real de la demanda y el cambio climático.}

\sol{85}{D}{\textbf{Tema:} Causales financieras de ingreso al RAT de OTASS. \textbf{Explicación:} El Decreto Legislativo 1280 establece causales objetivas de ingreso al RAT: insolvencia técnica y financiera (patrimonio neto menor a cero o pérdidas de más de dos tercios del capital), quiebra operativa y grave incapacidad de gestión del prestador municipal.}

\sol{86}{B}{\textbf{Tema:} Obras por Impuestos (OxI) en saneamiento. \textbf{Explicación:} OxI agiliza la inversión pública: una empresa privada adelanta el capital y ejecuta la obra de agua y alcantarillado, recibiendo del MEF Certificados de Inversión Pública Regional y Local (CIPRL) aplicables al pago de su impuesto a la renta.}

\sol{87}{A}{\textbf{Tema:} Contingencias ambientales y pasivos ocultos. \textbf{Explicación:} No provisionar contingencias por descargas de aguas residuales sin tratamiento (sancionables por OEFA y ANA) infla artificialmente el patrimonio de la EPS, engañando a auditores y retrasando la inversión necesaria en remediación ambiental.}

\sol{88}{A}{\textbf{Tema:} Evaluación por CAE en micromedición. \textbf{Explicación:} La comparación entre tecnologías de medición con distinta durabilidad y costos exige el Costo Anual Equivalente neto del beneficio de reducción de ANF aparente. La tecnología con menor costo neto anualizado maximiza la rentabilidad operativa de la EPS.}

\sol{89}{C}{\textbf{Tema:} Fundamento financiero del periodo regulatorio quinquenal. \textbf{Explicación:} Las inversiones en infraestructura hídrica tardan varios años en ejecutarse y madurar. Un quinquenio brinda certidumbre regulatoria a la empresa para invertir y capturar ganancias de eficiencia antes de que el regulador reajuste la tarifa base en el siguiente ciclo.}

\sol{90}{D}{\textbf{Tema:} Equilibrio económico-financiero en concesiones. \textbf{Explicación:} Las cláusulas de equilibrio económico-financiero blindan la bancabilidad de los contratos a largo plazo: si el Estado altera ex-post las condiciones regulatorias o tributarias de forma sustancial (hecho del príncipe), el contrato prevé ajustes en la tarifa o plazos para restituir el VAN original pactado.}

\subsection*{Bloque D. Ampliación: Teoría y Aplicaciones Prácticas (91 a 180)}

\sol{91}{B}{Por el límite notable $\lim_{m\to\infty}(1 + r/m)^m = e^r$, la tasa efectiva anual bajo capitalización continua es exactamente $TEA = e^r - 1$.}
\sol{92}{C}{Por Fisher, $1 + r_{ex-post} = \frac{1 + i}{1 + \pi_{efectiva}}$. Si $\pi_{efectiva} > \pi^e$, la tasa real efectiva cae y el deudor paga con moneda de menor poder adquisitivo.}
\sol{93}{D}{El crecimiento solo crea valor ($PVGO > 0$) cuando los proyectos reinvertidos generan $VAN > 0$, es decir, $ROE > k_e$. Si $ROE = k_e$, el $PVGO = 0$.}
\sol{94}{A}{La fórmula estándar de anualidad vencida coloca el valor presente un periodo antes del primer pago (en $t=k$); luego se trae a $t=0$ dividiendo entre $(1+r)^k$.}
\sol{95}{C}{La teoría de preferencia por la liquidez añade una prima positiva creciente con el plazo para compensar el mayor riesgo de precio de bonos largos, inclinando la curva hacia arriba.}
\sol{96}{D}{Sin arbitraje, invertir a 2 años a la tasa $s_2$ debe rendir igual que invertir a 1 año a $s_1$ y pactar la reinversión del año 2 a $f_{1,2}$: $(1+s_2)^2 = (1+s_1)(1+f_{1,2})$.}
\sol{97}{A}{En el sistema francés la cuota $C$ es constante; al reducirse el saldo deudor, los intereses bajan y la amortización sube. En el alemán la amortización es fija, bajando el saldo más rápido (menos intereses totales).}
\sol{98}{B}{A mayor $m$, los intereses se capitalizan antes y generan intereses sobre intereses, elevando la TEA de forma monótona y cóncava hacia su cota superior $e^{TNA}-1$.}
\sol{99}{D}{La utilidad neta se rige por el devengado (registra facturación aun si no se cobra) y no refleja desembolsos de capital (CAPEX) ni amortizaciones de principal de deuda.}
\sol{100}{A}{El FCL pertenece a todos los proveedores de capital (acreedores y accionistas), por lo que se descuenta al WACC; el FCA es residual para accionistas tras deuda y se descuenta a $k_e$.}
\sol{101}{B}{Aumentar el CTN (más inventarios o cuentas por cobrar netas) inmoviliza caja ($-\Delta CTN$ en el FCL); al finalizar el proyecto ese capital de trabajo se liquida y se recupera como flujo positivo.}
\sol{102}{C}{Al deducirse antes de impuestos y luego sumarse de vuelta al NOPAT, la depreciación incrementa el flujo de caja libre en exactamente $T \cdot Dep$ (escudo fiscal de la depreciación).}
\sol{103}{A}{Como $EVA = (ROIC - WACC) \times \text{Capital Invertido}$, si $ROIC < WACC$ la empresa no cubre el costo de oportunidad de su capital y destruye valor económico ($EVA < 0$).}
\sol{104}{B}{La identidad DuPont de 5 pasos expresa $ROE = \frac{UN}{EBT} \times \frac{EBT}{EBIT} \times \frac{EBIT}{\text{Ventas}} \times \frac{\text{Ventas}}{\text{Activos}} \times \frac{\text{Activos}}{\text{Patrimonio}}$.}
\sol{105}{C}{Por la relación $ROE = ROIC + (ROIC - k_d(1-T))\frac{D}{E}$, endeudarse eleva el $ROE$ si y solo si los activos rinden más que el costo neto de la deuda ($ROIC > k_d(1-T)$).}
\sol{106}{D}{$GAO = \frac{\%\Delta EBIT}{\%\Delta Q}$ (riesgo por costos fijos operativos) y $GAF = \frac{\%\Delta UPA}{\%\Delta EBIT}$ (riesgo por gastos financieros fijos); además $GAT = GAO \times GAF$.}
\sol{107}{B}{El uso de un activo existente tiene un costo de oportunidad igual a su mejor uso alternativo (valor de mercado neto de impuestos hoy), a diferencia de los costos hundidos irrecuperables.}
\sol{108}{C}{Con dos cambios de signo ($-, +, \dots, +, -$), la ecuación $VAN(r)=0$ puede tener hasta 2 raíces reales positivas (múltiples TIR). El VAN y la TIRM evitan esta ambigüedad.}
\sol{109}{D}{Como $VAN_A(r) - VAN_B(r) = VAN_{A-B}(r) = 0$, la tasa de Fisher $r_F$ es exactamente la TIR del flujo incremental $(FC_A - FC_B)$.}
\sol{110}{A}{Con presupuesto limitado en un periodo y proyectos divisibles, ordenar por $IR$ (o razón $VAN/I_0$) maximiza el VAN total por unidad de capital escaso (si son indivisibles, se usa programación entera).}
\sol{111}{C}{Un activo que dura 4 años deberá reemplazarse antes que uno de 6 años; el CAE anualiza el VAC a lo largo de la vida útil (equivalente a igualar horizontes mediante cadena de reemplazos).}
\sol{112}{D}{Se reemplaza cuando el costo incremental de operar el activo viejo un año más (O\&M + pérdida de valor de rescate + costo de oportunidad del rescate) excede el CAE del activo nuevo.}
\sol{113}{A}{Aunque descuenta los flujos al costo de capital, el Payback descontado descarta todo flujo posterior al periodo de corte, pudiendo rechazar proyectos de gran VAN con flujos tardíos.}
\sol{114}{B}{El análisis de sensibilidad varía una variable a la vez (*ceteris paribus*); Monte Carlo modela distribuciones conjuntas y correlaciones para obtener la distribución de probabilidad completa del VAN.}
\sol{115}{D}{El VAN expandido equivale a $VAN_{estático} + \text{Valor de Opciones Reales}$; el VAN tradicional omite la capacidad de diferir, expandir o abandonar conforme se resuelve la incertidumbre.}
\sol{116}{A}{Como $\sigma_p^2 = w_1^2\sigma_1^2 + w_2^2\sigma_2^2 + 2w_1 w_2 \rho_{12}\sigma_1\sigma_2$, para cualquier $\rho_{12} < 1$ se cumple $\sigma_p < w_1\sigma_1 + w_2\sigma_2$, existiendo ganancia de diversificación.}
\sol{117}{B}{La decisión de inversión (elegir el portafolio tangente de máximo Sharpe) se separa de la decisión de financiamiento/riesgo (combinar ese portafolio con el activo libre de riesgo sobre la CML).}
\sol{118}{C}{La CML relaciona $E[R_p]$ con $\sigma_p$ (pendiente = ratio de Sharpe del mercado) solo para portafolios eficientes; la SML relaciona $E[R_i]$ con $\beta_i$ (pendiente = $E[R_m]-r_f$) para todo activo en equilibrio.}
\sol{119}{A}{$\alpha_i = E[R_i] - [r_f + \beta_i(E[R_m]-r_f)] > 0$ significa que el activo promete más retorno del exigido por su beta: está subvaluado (barato), atrayendo compras que elevan su precio y bajan su retorno hacia la SML.}
\sol{120}{B}{Sharpe divide el exceso de retorno por el riesgo total ($\sigma_p$); Treynor lo divide por el riesgo sistemático ($\beta_p$), apropiado cuando el riesgo idiosincrático ya está diversificado.}
\sol{121}{C}{Hamada separa riesgo operativo ($\beta_U$) del financiero: primero se desapalanca cada comparable ($\beta_U = \frac{\beta_L}{1+(1-T)D/E}$) y luego se reapalanca con la $D/E$ y $T$ regulatorias.}
\sol{122}{D}{Sin impuestos, $k_e = k_u + (k_u - k_d)(D/E)$. Sustituyendo en $WACC = \frac{E}{V}k_e + \frac{D}{V}k_d$, los efectos se cancelan exactamente y $WACC = k_u$ permanece constante.}
\sol{123}{B}{En el modelo Trade-off, $V_L = V_U + VP(\text{Escudo Fiscal}) - VP(\text{Costos de Quiebra})$; el óptimo interior $D^*$ iguala el beneficio marginal tributario al costo marginal de insolvencia.}
\sol{124}{C}{Por selección adversa, emitir acciones señala al mercado que la gerencia las cree sobrevaluadas (haciendo caer el precio); por ello se prefieren fondos internos, luego deuda y al final acciones.}
\sol{125}{D}{Con responsabilidad limitada, el patrimonio actúa como una opción *call* sobre los activos: cerca de la quiebra, mayor volatilidad aumenta el valor del patrimonio a expensas de los acreedores.}
\sol{126}{A}{Sin fricciones, si un accionista desea efectivo cuando la firma no paga dividendo, vende una fracción de sus acciones ('dividendo casero'), dejando intacta su riqueza total.}
\sol{127}{C}{El WACC constante asume que $D/V$ se rebalancea continuamente. El APV ($VAN_{desapalancado} + VP(\text{efectos financieros})$) maneja con exactitud calendarios de deuda fijos, subsidios o *Project Finance*.}
\sol{128}{D}{Cuando la tasa exigida por el mercado supera la tasa cupón ($YTM > c$), el bono debe ofrecer una ganancia de capital para compensar, cotizando por debajo de su valor nominal (bajo par).}
\sol{129}{A}{La duración es el promedio ponderado temporal de los flujos: cupones altos adelantan el peso temporal (baja $D$); sin cupones (cupón cero), todo el flujo ocurre en $n$, por lo que $D = n$.}
\sol{130}{B}{La recta tangente (Duración Modificada) yace siempre por debajo de la curva convexa precio-YTM; por tanto, siempre subestima el precio real (subestima la ganancia si baja $y$ y sobreestima la pérdida si sube $y$).}
\sol{131}{D}{Cuando la duración de la cartera iguala el horizonte $H$, el efecto precio y el efecto reinversión ante un cambio pequeño y paralelo en las tasas se compensan mutuamente en $t = H$.}
\sol{132}{A}{En $t=T$, ambos portafolios pagan exactamente $\max(S_T, K)$. Por ley de un solo precio (no arbitraje), deben costar exactamente lo mismo hoy en $t=0$.}
\sol{133}{B}{Las opciones tienen pagos asimétricos (pérdida acotada a la prima y ganancia potencial ilimitada); mayor volatilidad aumenta la probabilidad de pagos favorables, elevando tanto *call* como *put* ($\text{Vega} > 0$).}
\sol{134}{C}{Un *Forward* es una obligación simétrica sin prima inicial; una *Call* otorga el derecho (no obligación) de comprar a $K$, protegiendo contra alzas del dólar y permitiendo beneficiarse si baja, a cambio de una prima.}
\sol{135}{A}{Esperar para invertir equivale a una *call* (derecho a pagar $I$ para adquirir el proyecto); abandonar por el valor de salvamento $S$ equivale a una *put* (derecho a vender el proyecto al precio $S$).}
\sol{136}{B}{SUNASS fija la tarifa media igualando el valor presente de ingresos proyectados al valor presente de costos eficientes (O\&M, inversiones financiadas por tarifa e impuestos) descontados al WACC real.}
\sol{137}{C}{Al no cotizar en bolsa y tener estructuras financieras distorsionadas, SUNASS emplea betas de empresas comparables (sector agua/utilities) ajustados por Hamada a una estructura de capital objetivo eficiente.}
\sol{138}{D}{El VNR valora la infraestructura eficiente requerida hoy a precios actuales de reposición, dando señales económicas correctas frente a la distorsión inflacionaria de los costos históricos.}
\sol{139}{B}{Los activos donados no implican costo de oportunidad de capital para la EPS (no se remuneran con WACC para no cobrar dos veces al ciudadano), pero sí requieren cubrir O\&M y futura reposición.}
\sol{140}{C}{El PMO planifica infraestructura y oferta-demanda en un horizonte de largo plazo (30 años), sirviendo de base para determinar la fórmula tarifaria, estructura tarifaria y metas de gestión cada 5 años.}
\sol{141}{D}{Si la tasa permitida supera el costo de capital, la regulación por tasa de retorno incentiva sobreinversión en capital (Averch-Johnson) y traslado de costos; *Price Cap* incentiva eficiencia en costos.}
\sol{142}{A}{La evaluación social mide el impacto neto en el bienestar del país: usa precios sombra (sociales), incorpora externalidades (ej. salud pública) y descuenta los flujos a la TSD fijada por el MEF.}
\sol{143}{C}{Desde la perspectiva del país en su conjunto, un impuesto o subsidio es una transferencia entre privados y el Estado; el costo económico real corresponde a los recursos netos de distorsiones fiscales.}
\sol{144}{D}{Donde hay subempleo, el costo de oportunidad social de la mano de obra no calificada es inferior al salario de mercado ($FC < 1$). Los impuestos/aranceles hacen que el valor social de una divisa supere al tipo de cambio de mercado ($FC > 1$).}
\sol{145}{A}{El beneficio social de pasar de cisterna (alto costo) a red pública (menor costo) es el área bajo la curva de demanda: recursos liberados del consumo antiguo más excedente por consumo incremental.}
\sol{146}{B}{Invierte.pe usa Costo-Beneficio cuando los beneficios son monetizables (agua potable); usa Costo-Eficacia (menor costo social por beneficiario o $m^3$ tratado) cuando monetizar beneficios resulta complejo (alcantarillado/PTAR).}
\sol{147}{D}{Una TSD de $8\%$ castiga drásticamente beneficios intergeneracionales (a 50 años, $1.08^{-50} \approx 0.02$). Por equidad intergeneracional e incertidumbre de largo plazo, el MEF establece tasas sociales menores/decrecientes para proyectos ambientales.}
\sol{148}{A}{El *Value for Money* compara el costo total del ciclo de vida (inversión, O\&M y riesgos transferidos) entre APP y obra pública tradicional (Comparador Público-Privado), buscando eficiencia neta.}
\sol{149}{B}{Transferir riesgos incontrolables al privado encarece innecesariamente el contrato (prima de riesgo excesiva). La regla óptima asigna cada riesgo a quien mejor puede controlarlo al menor costo.}
\sol{150}{C}{En esquemas de pago por disponibilidad, el concesionario recibe el RPI/RPMO (respaldado usualmente por fideicomiso) sujeto a cumplir los indicadores de disponibilidad y calidad (descuentos por incumplimiento).}
\sol{151}{A}{En OxI, el privado adelanta el financiamiento y ejecución de un proyecto público viable (Invierte.pe) y recibe CIPRL/CIPGN que aplica contra su Impuesto a la Renta, descontándose luego del canon/recursos de la entidad.}
\sol{152}{B}{Un *Working Ratio* $< 1$ (idealmente $< 0.70$) implica que los ingresos operativos cubren con holgura los costos operativos en efectivo, dejando excedente (margen EBITDA positivo) para servicio de deuda e inversiones.}
\sol{153}{C}{La intangibilidad obliga a depositar los porcentajes tarifarios asignados a MERESE/GRD en cuentas exclusivas; su desvío a gasto corriente está prohibido y sujeto a sanción por SUNASS.}
\sol{154}{D}{$TEA = \left(1 + \frac{0.12}{4}\right)^4 - 1 = (1.03)^4 - 1 = 1.125509 - 1 = 12.55\%$.}
\sol{155}{B}{Bajo capitalización continua: $VF = VP \cdot e^{rt} = 1{,}000 \cdot e^{0.08 \times 5} = 1{,}000 \cdot e^{0.40} \approx S/\ 1{,}491.82$.}
\sol{156}{C}{Por Fisher exacto: $1 + r = \frac{1 + i}{1 + \pi} = \frac{1.1025}{1.05} = 1.05 \Rightarrow r = 5.00\%$.}
\sol{157}{D}{$VP = 1{,}000 \times \frac{1 - (1.10)^{-3}}{0.10} = 1{,}000 \times \frac{1 - 0.751315}{0.10} = S/\ 2{,}486.85$.}
\sol{158}{A}{Primero hallamos $D_1 = D_0(1+g) = 4.00(1.05) = S/\ 4.20$. Luego, $P_0 = \frac{D_1}{k_e - g} = \frac{4.20}{0.13 - 0.05} = \frac{4.20}{0.08} = S/\ 52.50$.}
\sol{159}{C}{$g = ROE \times b = 0.15 \times 0.40 = 6.0\%$. El dividendo esperado es $D_1 = 0.60 \times 10 = S/\ 6.00$, y $P_0 = \frac{6.00}{0.12 - 0.06} = S/\ 100.00$.}
\sol{160}{D}{$CTN = AC - PC = 450 - 300 = S/\ 150$ millones; Razón Corriente $= \frac{AC}{PC} = \frac{450}{300} = 1.50$.}
\sol{161}{A}{La Prueba Ácida excluye los inventarios: $\frac{AC - \text{Inventarios}}{PC} = \frac{320 - 80}{200} = \frac{240}{200} = 1.20$.}
\sol{162}{B}{$ROE = \text{Margen Neto} \times \text{Rotación} \times \text{Multiplicador} = 0.08 \times 1.25 \times 1.80 = 0.18 = 18.0\%$.}
\sol{163}{D}{$FCL = EBIT(1-T) + Dep - CAPEX - \Delta CTN = 200(0.70) + 50 - 80 - 20 = 140 + 50 - 100 = S/\ 90$ mil.}
\sol{164}{A}{$EVA = NOPAT - (WACC \times \text{Capital Invertido}) = 65 - (0.10 \times 500) = 65 - 50 = S/\ 15$ millones.}
\sol{165}{B}{$VAN = -1{,}000 + \frac{550}{1.10} + \frac{726}{(1.10)^2} = -1{,}000 + 500 + 600 = S/\ 100.00$.}
\sol{166}{C}{El Índice de Rentabilidad es $IR = \frac{VP(\text{Flujos})}{I_0} = \frac{1{,}320}{1{,}000} = 1.32$ (como $IR > 1$, el proyecto crea valor).}
\sol{167}{A}{Resolviendo $-100 + \frac{60}{1+r} + \frac{72}{(1+r)^2} = 0$ con $x = 1+r$: $25x^2 - 15x - 18 = (5x-6)(5x+3) = 0 \Rightarrow x = 1.20 \Rightarrow TIR = 20.0\%$.}
\sol{168}{B}{Igualando $VAN_A = VAN_B$ (flujo incremental $B-A$: $0, -110, +121$): $-\frac{110}{1+r_F} + \frac{121}{(1+r_F)^2} = 0 \Rightarrow 1 + r_F = \frac{121}{110} = 1.10 \Rightarrow r_F = 10.0\%$.}
\sol{169}{C}{$CAE = VAC \times \frac{r}{1 - (1+r)^{-n}} = 1{,}000 \times \frac{0.10}{1 - 0.7513} = \frac{100}{0.2487} \approx S/\ 402.11$.}
\sol{170}{D}{Los valores presentes son $VP_1 = 500$, $VP_2 = 400$ (acumulado $900$) y $VP_3 = 200$. Faltan $S/\ 100$ en el año 3: $2 + \frac{100}{200} = 2.50$ años.}
\sol{171}{B}{Aplicando el CAPM emergente: $k_e = r_f + \beta_L(E[R_m] - r_f) + RP = 4.0\% + 0.80(6.0\%) + 2.0\% = 10.80\%$.}
\sol{172}{C}{$\beta_L = \beta_U\left[1 + (1 - T)\frac{D}{E}\right] = 0.60[1 + (1 - 0.30)(0.50)] = 0.60(1.35) = 0.81$.}
\sol{173}{D}{$WACC = \frac{E}{V}k_e + \frac{D}{V}k_d(1-T) = 0.60(12.0\%) + 0.40(8.0\%)(0.75) = 7.20\% + 2.40\% = 9.60\%$.}
\sol{174}{A}{Según MM I con impuestos corporativos y deuda perpetua: $V_L = V_U + T \cdot D = 800 + 0.30(200) = S/\ 860$ millones.}
\sol{175}{C}{Por MM II sin impuestos: $k_e = k_u + (k_u - k_d)\frac{D}{E} = 11.0\% + (11.0\% - 6.0\%)(0.80) = 11.0\% + 4.0\% = 15.0\%$.}
\sol{176}{D}{$P_B = \frac{80}{1.10} + \frac{1{,}080}{(1.10)^2} = 72.73 + 892.56 = S/\ 965.29$ (cotiza bajo par porque $c < YTM$).}
\sol{177}{A}{Primero, $D_{mod} = \frac{5.40}{1.08} = 5.00$. Luego, $\frac{\Delta P}{P} \approx -D_{mod}\Delta y = -5.00(0.005) = -0.025 = -2.50\%$.}
\sol{178}{B}{Por paridad: $P = C + \frac{K}{1+r} - S_0 = 12.00 + \frac{110.25}{1.05} - 100.00 = 12.00 + 105.00 - 100.00 = S/\ 17.00$.}
\sol{179}{D}{Retorno CAPM esperado: $4.0\% + 1.25(10.0\% - 4.0\%) = 11.5\%$, luego $\alpha_p = 14.0\% - 11.5\% = +2.5\%$. El Ratio de Sharpe es $SR = \frac{14.0\% - 4.0\%}{20.0\%} = 0.50$.}
\sol{180}{A}{Aplicando cada factor de corrección: $600(0.85) + 200(0.50) + 200(1.10) = 510 + 100 + 220 = S/\ 830$ mil.}


\end{document}
